\documentclass[11pt,a4paper]{report}
% \usepackage{polski}
\usepackage[utf8]{inputenc}
\usepackage[margin=3.5cm]{geometry}
\usepackage{graphicx}
\usepackage{amsmath}

%% ########################################################
\begin{document}
\thispagestyle{empty}
\begin{center}
\includegraphics[width=\textwidth]{logo_AGH.pdf}\\
{\bf{\sf WYDZIAŁ ELEKTROTECHNIKI, AUTOMATYKI, INFORMATYKI I INŻYNIERII BIOMEDYCZNEJ}}\\[5mm]
%% ======================================================
{\bf{\sf{KATEDRA AUTOMATYKI I ROBOTYKI}}}\\[14mm]
%% KATEDRA FIZYKI CIAŁA STAŁEGO
%% KATEDRA FIZYKI MATERII Skondensowanej
%% KATEDRA FIZYKI MEDYCZNEJ I BIOFIZYKI
%% KATEDRA INFORMATYKI STOSOWANEJ I FIZYKI KOMPUTEROWEJ
%% KATEDRA ODDZIAŁYWAŃ I DETEKCJI CZĄSTEK
%% KATEDRA ZASTOSOWAŃ FIZYKI JĄDROWEJ
%% ======================================================
{\sf{\huge Praca dyplomowa}}\\[12mm] 
%% Projekt dyplomowy = inżynierska
%% Praca dyplomowa = magisterska
%% ======================================================
{\sf{\Large Comparison of uncertainty approximation methods in machine learning and Bayesian methods\\[2mm] %% Niepotrzebne usunąć
*Polska wersja*\\[2mm] %% Jeżeli praca została napisana w języku innym niż język polski
 %% Jeżeli praca została napisana w języku innym niż język angielski
%% ======================================================
%% W przypadku pracy napisanej
%% - po polsku: dwa tytuły,
%% - po angielsku: dwa tytuły,
%% - po hiszpańsku/niemiecku/rosyjsku/etc: trzy tytuły.
%% ======================================================
}}\\[40mm]
\end{center}
{\sf
\begin{tabular}{ll}
Autor: & Michał Kocwa\\
Kierunek studiów: &	Automatyka i Robotyka\\
%% Fizyka Techniczna
%% Fizyka Medyczna
%% Informatyka Stosowana 
%% Mikro- i Nanotechnologie w biofizyce
%% Nanoinżynieria Materiałów
%% Computer Physics
Opiekun pracy: & dr Waldemar Bauer \\
\end{tabular}
}\\[10mm]
\begin{center}
{\sf Kraków, 2026}
\end{center}

%% ########################################################
\newpage

\tableofcontents
\newpage
% Rozdziały
\chapter{Introduction}

\section{Motivation and background}

\section{Aims and scope}



\chapter{Theoretical Background of Uncertainty in Machine Learning}


\section{Types of uncertainty}
\label{sec:types-of-uncertainty}

The value obtained from the predictive model rarely reproduces the observed value exactly. The difference between the predicted and the ground truth values has two distinct sources, and separating them is the starting point for this entire work. Considering two test cases for which a regression model returns prediction intervals of the same width, in the first case, the training data are dense but scattered around the trend, while in the second one, the model extrapolates far beyond the range of the training data, the density being lower. A single interval width conceals this difference. The appropriate response applied to both cases fundamentally differs: collecting more data helps in the second case, but not in the first, where the amount of available data points exceeds the second one. Therefore, uncertainty must be classified by its source, and not merely by its magnitude.

To precisely capture both sources, let the data be generated according to the model
%
\begin{equation}
    y = f(x) + \varepsilon ,
    \label{eq:generative-model}
\end{equation}
%

where $f$ is the unknown target function and $\varepsilon$ is the noise component that appears in the observations. Uncertainty is introduced through two independent channels: through the noise $\varepsilon$, which no model can eliminate, and through the unknown function $f$, which the model must estimate from a finite amount of data. These two channels give rise to the two types of uncertainty discussed throughout this chapter~\cite{hullermeier2021, he2025}.
\begin{itemize}
    \item The first type is \emph{aleatoric uncertainty}, which stems from the randomness inherent in the data generation process. It cannot be reduced by collecting more observations. 
    \item The second is \emph{epistemic uncertainty}, which is due to the limited knowledge of the model and can be reduced with sufficient data~\cite{kendall2017}. 
\end{itemize}

In the literature, one may come across a notation that equates \emph{data uncertainty} and \emph{model uncertainty} as synonyms for the aleatoric and epistemic types, respectively~\cite{he2025, gawlikowski2023}. To avoid ambiguity, this work adopts the terms aleatoric and epistemic as primary, while reserving the term \emph{model uncertainty} for a subtype of epistemic uncertainty in the sense introduced in Subsection~\ref{subsec:epistemic}.

It should be noted that the boundary between the two types is not a property of the world, but rather of the chosen model. A source of variability that is aleatoric for one model can become epistemic for a richer model equipped with an additional feature~\cite{derkiureghian2009}. The distinction is therefore a modeling decision, made to indicate which portion of the total uncertainty can in practice be reduced. This pragmatic interpretation structures the subsequent discussion.

% ---------------------------------------------------------------------
\subsection{Aleatoric uncertainty}
\label{subsec:aleatoric}

The term \emph{aleatoric} comes from the Latin \emph{alea}, dice, and refers to uncertainty rooted in intrinsic randomness. Aleatoric uncertainty captures the noise inherent in the observations~\cite{kendall2017}. Because this noise belongs to the data generation process itself, it persists regardless of the number of collected examples. For this reason, it is also referred to as irreducible or statistical uncertainty. The altitude uncertainty has several sources. Sensor imprecision and measurement error distort the recorded values. Quantization discretizes a continuous quantity. Relevant features may be missing from the input vector, causing seemingly identical inputs to correspond to different target values. In classification, this same effect takes the form of class overlap, where a single region of the feature space is shared by more than one label. Another frequently cited source is the disagreement among data annotators. Two radiologists may contour the same lesion differently and no increase in the number of patients will eliminate this disagreement. Such cases show that aleatoric uncertainty reflects the data, rather than a deficiency of the model.In the regression formulation of Equation~\eqref{eq:generative-model}, the aleatoric uncertainty is carried by the noise term $\varepsilon$. Two cases are usually distinguished. In the homoscedastic case, the noise has a constant variance,
%
\begin{equation}
    \varepsilon \sim \mathcal{N}\!\left(0, \sigma^2\right),
    \label{eq:homoscedastic}
\end{equation}
%
so each observation is equally noisy. In the \emph{heteroscedastic} case, the variance depends on the input,
%
\begin{equation}
    \varepsilon(x) \sim \mathcal{N}\!\left(0, \sigma^2(x)\right),
    \qquad
    p(y \mid x) = \mathcal{N}\!\left(y;\, f(x),\, \sigma^2(x)\right),
    \label{eq:heteroscedastic}
\end{equation}
%
so some regions of the input space are inherently noisier than others. The heteroscedastic model is more general of the two. Heteroscedastic aleatoric uncertainty can be estimated directly with a network with two output heads: one for the predictive mean $\hat{\mu}_\theta(x)$, and the other for the predictive variance $\hat{\sigma}^2_\theta(x)$. The parameters are trained by minimizing the Gaussian negative log-likelihood,
%
\begin{equation}
    \mathcal{L}(\theta)
    = \frac{1}{N} \sum_{i=1}^{N}
      \left[
        \frac{\bigl(y_i - \hat{\mu}_\theta(x_i)\bigr)^2}
             {2\,\hat{\sigma}^2_\theta(x_i)}
        + \frac{1}{2}\,\log \hat{\sigma}^2_\theta(x_i)
      \right] .
    \label{eq:gaussian-nll}
\end{equation}
%
This function has an intuitive interpretation, sometimes called \emph{learned attenuation}~\cite{kendall2017}. The first term is the squared error weighted by the inverse of the predicted variance, so the network can reduce the weight of a large error by using a high variance. The second term, $\tfrac{1}{2}\log \hat{\sigma}^2_\theta(x)$, penalizes this move and prevents the variance from growing without bound. Together, both terms force the network to report high variance only where the data truly justify it. A limiting case clarifies the connection to standard regression. If the variance is constant, the second term of Equation~\eqref{eq:gaussian-nll} becomes a constant, and the first reduces to a squared error. The objective function then simplifies to mean squared error (MSE). Standard regression with a squared-error loss function thus already assumes a homoscedastic Gaussian noise model, even if this assumption remains implicit. The variance head removes this assumption and allows the noise level to vary with input. In classification, the predictive distribution $p(y \mid x)$ is categorical, and the aleatoric uncertainty is quantified by its entropy,
%
\begin{equation}
    H\!\left[p(y \mid x)\right]
    = -\sum_{c=1}^{K} p(y = c \mid x)\,\log p(y = c \mid x) .
    \label{eq:entropy}
\end{equation}
%
Entropy is maximized when the classes are equally likely, which corresponds precisely to the situation of complete class overlap. It vanishes when all probability masses are assigned to a single class.

\begin{figure}[htbp] % poprawić miejsce wykresu/podrozdzialu !!
\centering
\includegraphics[width=\textwidth]{rodzial2_rys/img2_1.png}
\caption{Heteroscedastic aleatoric uncertainty evaluated on synthetic data (sine function with noise increasing with $x$). The predictive mean curve is shown with a $\pm 2\hat{\sigma}_\theta(x)$ band that widens as observation noise grows, alongside the training points. The band width reflects noise inherent to the data rather than a lack of observations.}
\label{fig:heteroscedastic-aleatoric}
\end{figure}


% ---------------------------------------------------------------------
\subsection{Epistemic uncertainty}
\label{subsec:epistemic}

The term \emph{epistemic} comes from the Greek \emph{episteme}, knowledge, and refers to uncertainty resulting from a lack thereof. Epistemic uncertainty captures the uncertainty about the model and can be eliminated with sufficient data~\cite{kendall2017}. If training data are dense, the model is well constrained and epistemic uncertainty is low, but in the case where there training set is lacking data, many different models remain consistent with the observations and epistemic uncertainty is high.A single point estimation of parameters cannot express this uncertainty. In an overparameterized network, many parameter sets fit the training data almost equally well. These sets agree in regions covered by the data and diverge outside them. A standard network returns the prediction of a single such set and does not determine that other, equally plausible sets would provide a different response. This is the mechanism behind the overconfidence of deterministic networks~\cite{he2025}.The Bayesian approach replaces a single estimate with a distribution over parameters. The prior distribution $p(\theta)$ encodes beliefs about the parameters before seeing any data. After observing the data set $\mathcal{D} = \{(x_i, y_i)\}_{i=1}^{N}$, the posterior distribution follows from Bayes' theorem,
%
\begin{equation}
    p(\theta \mid \mathcal{D})
    = \frac{p(\mathcal{D} \mid \theta)\, p(\theta)}{p(\mathcal{D})} ,
    \label{eq:bayes}
\end{equation}
%
where $p(\mathcal{D} \mid \theta)$ is the likelihood and $p(\mathcal{D})$ is the normalization constant. The prediction at a new point $x^\ast$ is then obtained by marginalizing over the posterior distribution,
%
\begin{equation}
    p(y^\ast \mid x^\ast, \mathcal{D})
    = \int p(y^\ast \mid x^\ast, \theta)\,
           p(\theta \mid \mathcal{D})\; \mathrm{d}\theta .
    \label{eq:predictive}
\end{equation}
%
Equation~\eqref{eq:predictive} averages the predictions in every set of parameters, weighting each by its posterior likelihood. The spread of this predictive distribution is a measure of epistemic uncertainty. Where plausible sets agree, the spread is small; where they diverge, they are large. Exact computation of the integral is intractable for neural networks, and the methods compared in Chapter~3 differ mainly in how they approximate it. Epistemic uncertainty itself admits a finer breakdown~\cite{hullermeier2021}. Let $\mathcal{F}$ denote the space of all functions mapping inputs to target values, and $\mathcal{H} \subseteq \mathcal{F}$—a smaller space realizable by the chosen model class. Three functions are of interest: the optimum $f^\ast$ in the space $\mathcal{F}$, the optimum $h^\ast$ in the space $\mathcal{H}$, and the function $\hat{h}$ actually learned from finite data. The differences between them separate two contributions. The difference between $f^\ast$ and $h^\ast$ is \emph{the uncertainty} of the model, which originates from restricting attention to $\mathcal{H}$ and corresponds to the bias of the model class. The difference between $h^\ast$ and $\hat{h}$ is \emph{the approximation uncertainty}, which arises from learning on a finite sample and corresponds to variance.Both contributions are reduced by different means. The approximation uncertainty decreases as the training set grows, in the usual sense of a reducible quantity. Model uncertainty does not behave this way; it is reduced only by enlarging the model class, for example, by choosing a deeper architecture. For highly flexible model classes used in deep learning, the space $\mathcal{H}$ is sufficiently expressive for $h^\ast$ to lie close to $f^\ast$, and it is usually assumed that the model uncertainty component is negligible~\cite{hullermeier2021}. This work adopts this assumption and focuses on parameter uncertainty, which the Bayesian methods in Chapter~3 aim to capture.

\begin{figure}[htbp]
    \centering
    \includegraphics[width=\textwidth]{rodzial2_rys/img2_2.png}
    \caption{Epistemic uncertainty illustrated as function spread. Several function samples from the posterior distribution (e.g., Gaussian process or Monte Carlo dropout) evaluated on the synthetic sine dataset. The sampled curves overlap closely within the training range $[0, 6]$ and fan out significantly beyond it (in regions $[-2, 0]$ and $[6, 10]$), demonstrating the increase in epistemic uncertainty during extrapolation.}
    \label{fig:epistemic-uncertainty}
\end{figure}

\begin{figure}[htbp]
    \centering
    \includegraphics[width=0.7\textwidth]{rodzial2_rys/img2_3_1.png}
    \caption{Conceptual diagram of nested hypothesis spaces. The space $\mathcal{F}$ represents all functions, while $\mathcal{H}$ denotes functions realizable by the chosen model architecture. Points $f^*$, $h^*$, and $\hat{h}$ mark the true target function, the optimal model within $\mathcal{H}$, and the function actually learned from finite data, respectively. Adapted from Hüllermeier and Waegeman (2021), Fig.~4.}
    \label{fig:hypothesis-spaces}
\end{figure}

Both types of uncertainty are not competing descriptions, but rather additive components of total predictive uncertainty. In the predictive model of this chapter, they are separated by the law of total variance,
%
\begin{equation}
    \underbrace{\operatorname{Var}\!\left(y^\ast \mid x^\ast\right)}_{\text{total}}
    =
    \underbrace{\mathbb{E}_{p(\theta \mid \mathcal{D})}
      \!\left[\sigma^2_\theta(x^\ast)\right]}_{\text{aleatoric}}
    +
    \underbrace{\operatorname{Var}_{p(\theta \mid \mathcal{D})}
      \!\left(\mu_\theta(x^\ast)\right)}_{\text{epistemic}} .
    \label{eq:total-variance}
\end{equation}
%
The first term is the expected predicted noise variance and represents the aleatoric component. The second term is the variance of the predicted mean in the posterior distribution and represents the epistemic component. Equation~\eqref{eq:total-variance} provides an operational recipe used in Chapter~4 to separately report both contributions for each method.

\begin{figure}[htbp]
    \centering
    \includegraphics[width=\textwidth]{rodzial2_rys/img2_3.png}
    \caption{Predictive uncertainty decomposition showing two nested bands: the inner band represents the aleatoric component, while the outer band represents the total uncertainty. The aleatoric component dominates in the center of the training range, whereas epistemic uncertainty dominates at the boundaries, illustrating the law of total variance from Equation~\eqref{eq:total-variance}.}
    \label{fig:hypothesis-spaces}
\end{figure}
\section{Bayesian Inference}
\label{sec:bayesian-inference}

Section~\ref{subsec:epistemic} proved that epistemic uncertainty requires a
distribution over parameters rather than a single estimate. This section provides the
formalism that defines such a distribution and indicates why limitations prohibit it from being computed directly. All the mathematical methods used to the compared methods are collected in this section. Chapter~3 refers to methods shown below and does not repeat the derivations.

% ---------------------------------------------------------------------
\subsection{Bayes' Rule}
\label{subsec:bayes-rule}

Let $\mathcal{D} = \{(x_i, y_i)\}_{i=1}^{N}$ denote the training set of ordered pairs from the Cartesian product of sets of labels and feautures and $\theta$ --- the model parameters. Bayes' rule connects expected results prior to observing data with expected results after observing them,
%
\begin{equation}
    p(\theta \mid \mathcal{D})
    = \frac{p(\mathcal{D} \mid \theta)\, p(\theta)}{p(\mathcal{D})},
    \qquad
    p(\mathcal{D}) = \int p(\mathcal{D} \mid \theta)\, p(\theta)\; \mathrm{d}\theta .
    \label{eq:bayes-rule}
\end{equation}
%
The prior distribution $p(\theta)$ encodes knowledge of the parameters before seeing the
data. The likelihood $p(\mathcal{D} \mid \theta)$ describes how well a given set of
parameters explains the observations. The posterior distribution $p(\theta \mid \mathcal{D})$
is the result of inference. The denominator $p(\mathcal{D})$, called the evidence or
marginal likelihood, serves as a normalizing constant. All subsequent computational
difficulties arise from this single component.

% ---------------------------------------------------------------------
\subsection{Bayesian Inference vs. Point Estimation}
\label{subsec:mle-map}

Standard neural network training returns a single set of weights, that corresponds to maximum
likelihood estimation (MLE),
%
\begin{equation}
    \theta_{\mathrm{MLE}} = \arg\max_{\theta} \, \log p(\mathcal{D} \mid \theta) .
    \label{eq:mle}
\end{equation}
%
Incorporating a prior distribution leads to maximum a posteriori (MAP) estimation,
%
\begin{equation}
    \theta_{\mathrm{MAP}}
    = \arg\max_{\theta} \,
      \bigl[\log p(\mathcal{D} \mid \theta) + \log p(\theta)\bigr] .
    \label{eq:map}
\end{equation}
%
The marginal likelihood $p(\mathcal{D})$ does not depend on $\theta$ and vanishes in both optimization problems. A Gaussian prior $p(\theta) = \mathcal{N}(0, \sigma_p^2 I)$ yields the term
$-\lVert \theta \rVert^2 / (2\sigma_p^2)$ after taking its logarithm, which corresponds to the L2
regularization. Training with weight decay thus implements the estimation of MAP, although it is
usually not presented in this way~\cite{jospin2022}.

Both procedures return a point, not a distribution, thus reducing the posterior value to a single value, which leads to removal of information on how many other sets of parameters explained the data equally well. This information is precisely what constitutes epistemic uncertainty. Bayesian inference consists of abandoning this simplification and preserving the entire distribution $p(\theta \mid \mathcal{D})$. However, MAP estimation remains important for selected methods: for example the Laplace approximation in Section~3.4 builds a Gaussian distribution around $\theta_{\mathrm{MAP}}$, thus treating the point as a starting point rather than a final result.

% ---------------------------------------------------------------------
\subsection{Bayesian Neural Network}
\label{subsec:bnn-formalism}

Transferring the scheme described before to a neural network requires three major changes: 
\begin{itemize}
    \item The randominitialization of weights is replaced by an explicit prior distribution $p(\theta)$, usually Gaussian and independent for each weight.
    \item A single set of weights after training is replaced by a posterior distribution.
    \item Prediction ceases to be a forward pass through the network and becomes an average over this distribution.
\end{itemize}
A model built this way is called a Bayesian neural network (BNN)~\cite{jospin2022}.

The likelihood stems from the assumed noise model. For heteroscedastic regression from
Equation~\eqref{eq:heteroscedastic}, it takes the form
%
\begin{equation}
    p(\mathcal{D} \mid \theta)
    = \prod_{i=1}^{N}
      \mathcal{N}\!\left(y_i;\, \mu_\theta(x_i),\, \sigma^2_\theta(x_i)\right),
    \label{eq:likelihood-regression}
\end{equation}
%
where $\mu_\theta$ and $\sigma^2_\theta$ follow the noise model of
Section~\ref{subsec:aleatoric}. In the experiments of this chapter and the next,
$\sigma^2_\theta$ is a single learnt constant rather than a function of the
input; Equation~\eqref{eq:likelihood-regression} is written in its general form
because the Bayesian treatment that follows does not depend on that choice. Maximizing the logarithm of expression
\eqref{eq:likelihood-regression} is equivalent to minimizing the Gaussian negative
logarithmic likelihood of Equation~\eqref{eq:gaussian-nll}. The Bayesian formulation thus does
not replace the noise model, but builds upon it. A network introducing distribution over
weights also reports high uncertainty where training data was lacking, not just where it was
noisy.

% ---------------------------------------------------------------------
\subsection{Predictive Distribution}
\label{subsec:predictive-distribution}

The quantity of interest to the model user is not the posterior, but the
prediction at a new point $x^\ast$, which is obtained by marginalizing over the posterior
distribution,
%
\begin{equation}
    p(y^\ast \mid x^\ast, \mathcal{D})
    = \int p(y^\ast \mid x^\ast, \theta)\,
            p(\theta \mid \mathcal{D})\; \mathrm{d}\theta .
    \label{eq:predictive-2}
\end{equation}
%
Equation~\eqref{eq:predictive-2} is the central formula of this work. A Bayesian prediction
is the average of the predictions of all admissible models, weighted by their plausibility in
light of the data. Each of the five methods compared in Chapter~3 represents a different way
of approximating this integral, and it is the only reason that these methods differ.

The spread of the predictive distribution decomposes into two components introduced in
Section~\ref{sec:types-of-uncertainty}. The law of total variance yields
%
\begin{equation}
    \underbrace{\operatorname{Var}\!\left(y^\ast \mid x^\ast\right)}_{\text{total}}
    =
    \underbrace{\mathbb{E}_{p(\theta \mid \mathcal{D})}
      \!\left[\sigma^2_\theta(x^\ast)\right]}_{\text{aleatoric}}
    +
    \underbrace{\operatorname{Var}_{p(\theta \mid \mathcal{D})}
      \!\left(\mu_\theta(x^\ast)\right)}_{\text{epistemic}} .
    \label{eq:total-variance-2}
\end{equation}
%
The first term is the expected predicted noise variance. The second is the variance of the
predicted mean over the posterior and vanishes when the posterior collapses to a single
point. The estimation of a single point from Section~\ref{subsec:mle-map} therefore sets the second term to zero by assumption, rather than by observation~\cite{hullermeier2021}. Equation
\eqref{eq:total-variance-2} serves as the operational recipe used in Chapter~4: for each
method, both components are reported separately.

% ---------------------------------------------------------------------
\subsection{Intractability of the Posterior}
\label{subsec:intractability}

Equations~\eqref{eq:bayes-rule} and~\eqref{eq:predictive-2} are exact, but for neural
networks, they cannot be computed directly. The source of difficulty is the integral in the
denominator of expression \eqref{eq:bayes-rule}, arising from three independent reasons:

\begin{itemize}
    \item dimensionality --- integration takes place over a parameter space of size
          on the order of $10^5$ and larger, where numerical quadratures are useless;
    \item lack of conjugacy --- network non-linearities mean that the product of the
          likelihood and the prior does not belong to any family with a known analytical
          form;
    \item posterior geometry --- network symmetries, such as neuron permutations within a
          layer, replicate every solution, making the posterior highly
          multimodal~\cite{jospin2022}.
\end{itemize}

The consequence is practical - since $p(\theta \mid \mathcal{D})$ remains unknown,
integral~\eqref{eq:predictive-2} cannot be calculated directly, making approximation
necessary. Chapter~3 compares methods that perform such an approximation; this section
justifies why they are necessary in the first place.

% ---------------------------------------------------------------------
\subsection{Two Approximation Strategies}
\label{subsec:mcmc-vs-vi}

Approximate inference methods fall into two families. The first samples from the exact
posterior without determining it explicitly. These include Markov Chain Monte Carlo
(MCMC) methods, including Hamiltonian Monte Carlo. Their advantage is convergence to the
true posterior given a sufficiently long chain. The second family replaces the posterior
with a distribution of a known form from which sampling is cheap, converting inference into
an optimization problem.

MCMC methods were omitted in this work. Therefore, all methods compared in Chapter~3 belong to the second
family.

% ---------------------------------------------------------------------
\subsection{Variational Inference *Może do usunięcia*}
\label{subsec:variational-inference}

Variational inference (VI) converts the integration problem into an optimization problem by assumpting existence of family of distributions $\mathcal{Q}$ of a known form, and the element
closest to the true posterior is sought~\cite{blei2017} using Kullback--Leibler divergence to measure its closeness.
%
\begin{equation}
    q^\ast(\theta)
    = \arg\min_{q \in \mathcal{Q}}\;
      \mathrm{KL}\bigl(q(\theta) \,\|\, p(\theta \mid \mathcal{D})\bigr) .
    \label{eq:vi-objective}
\end{equation}
%
Problem~\eqref{eq:vi-objective} in this form still contains the unknown posterior.
Expanding the divergence, however, leads to the identity
%
\begin{equation}
    \log p(\mathcal{D})
    = \mathcal{L}(q) + \mathrm{KL}\bigl(q(\theta) \,\|\, p(\theta \mid \mathcal{D})\bigr),
    \qquad 
    \mathcal{L}(q)
    = \mathbb{E}_{q}\bigl[\log p(\mathcal{D} \mid \theta)\bigr]
      - \mathrm{KL}\bigl(q(\theta) \,\|\, p(\theta)\bigr) .
    \label{eq:elbo}
\end{equation}
%
The left-hand side does not depend on $q$, and the divergence is non-negative, thus, the
quantity $\mathcal{L}(q)$ constitutes a lower bound on the evidence and is called the
evidence lower bound (ELBO). Minimizing the divergence is equivalent to maximizing the
ELBO, and the latter does not require knowledge of the posterior. Both terms of expression
$\mathcal{L}(q)$ have a clear interpretation: the first rewards data fitness, while the second
penalizes divergence from the prior distribution.

In practice, the family $\mathcal{Q}$ must be narrow enough to be manageable. The
most common choice is the mean-field approximation, in which the parameters are mutually
independent,
%
\begin{equation}
    q_\phi(\theta) = \prod_{j=1}^{M} q_{\phi_j}(\theta_j),
    \qquad
    q_{\phi_j}(\theta_j) = \mathcal{N}(\theta_j;\, \mu_j,\, \sigma_j^2) .
    \label{eq:mean-field}
\end{equation}
%
The expectation gradient in Equation~\eqref{eq:elbo} cannot be calculated directly
because the distribution that being averaged over itself depends on the parameters that
are optimized. The solution is the reparameterization: a sample is written as a
deterministic function of the parameters and a random variable with a distribution
independent of them,
%
\begin{equation}
    \theta = \mu + \sigma \odot \varepsilon,
    \qquad \varepsilon \sim \mathcal{N}(0, I) .
    \label{eq:reparameterization}
\end{equation}
%
In this way, randomness is transferred to $\varepsilon$, and the gradient with respect to
$\mu$ and $\sigma$ passes through standard backpropagation. This procedure makes VI
applicable to neural networks and forms the basis of the Bayes by Backprop algorithm in
Section~3.2.

A limitation of this strategy is fact that the solution $q^\ast$ is the best
approximation \emph{within the assumed family}, not an approximation with controlled error.
The direction of the divergence used in Equation~\eqref{eq:vi-objective} favors
distributions concentrated around a single mode, causing VI to systematically underestimate
the posterior variance~\cite{blei2017}. The mean-field approximation additionally ignores
correlations between weights. In a work dedicated to uncertainty estimation, this has direct
significance: an underestimated posterior spread translates to an underestimated epistemic
term in Equation~\eqref{eq:total-variance-2}. This observation recurs in Chapter~4 during the
interpretation of the PICP metric.

%\section{Types of uncertainty}
%\subsection{Aleatoric uncertainty}
%\subsection{Epistemic uncertainty}

%R\section{Bayesian inference}



\chapter{Uncertainty estimation methods}

\section{Gaussian processes}
\label{sec:gaussian-processes}

Section~\ref{subsec:intractability} showed that the posterior over the weights of a neural
network cannot be obtained in closed form. Gaussian processes circumvent this obstacle by
changing the object of inference - instead of a distribution over a finite set of parameters,
a distribution is placed directly over functions. For a Gaussian likelihood, both the
posterior and the predictive distribution of Equation~\eqref{eq:predictive-2} then admit exact
analytical solutions. Therefore, in this work, a Gaussian process (GP) serves as a reference
point, being the only one of the five compared methods whose uncertainty estimates do not involve an approximation of the posterior~\cite{rasmussen2006}.

% ---------------------------------------------------------------------
\subsection{A distribution over functions}
\label{subsec:gp-definition}

A Gaussian process is a collection of random variables, any finite number of which have a
joint Gaussian distribution~\cite{rasmussen2006}. 
\begin{equation}
    f(x) \sim \mathcal{GP}\bigl(m(x),\, k(x, x')\bigr),
\end{equation}
It is fully specified by a mean function
$m(x)$,
%
\begin{equation}
    m(x) = \mathbb{E}\left[f(x)\right],
    \label{eq:gp-definition}
\end{equation}
%
and a covariance function $k(x, x')$.
\begin{equation}
    k(x, x') = \mathbb{E}\bigl[(f(x) - m(x))(f(x') - m(x'))\bigr] .
    \label{eq:gp-covariance}
\end{equation}
%
The definition refers to an infinite-dimensional object, but it is used only through its
finite-dimensional marginals. For any finite set of inputs $X = \{x_1, \dots, x_N\}$, the
vector of function values follows a multivariate normal distribution,
%
\begin{equation}
    \mathbf{f} = \bigl[f(x_1), \dots, f(x_N)\bigr]^{\!\top}
    \sim \mathcal{N}\!\left(\mathbf{m},\, K\right),
    \qquad
    K_{ij} = k(x_i, x_j) .
    \label{eq:gp-marginal}
\end{equation}
%
This property is what makes computation possible because inference is carried out only at the points that actually occur in the training set and in the test set, and the remainder of the function is never represented explicitly.

The mean function is set to zero in the experiments in this work. The assumption is not
restrictive, because the targets are centered before training, thus all prior knowledge about the
shape of the function is carried by the covariance function alone.

% ---------------------------------------------------------------------
\subsection{Covariance functions}
\label{subsec:gp-kernel}

The covariance function, also called the kernel, determines which functions are plausible
before any data is seen. It states how strongly the values at two inputs are coupled to each other, and consequently how smooth and how variable the sampled functions are. 

The squared exponential kernel, also known as the radial basis function (RBF) kernel, is the most common choice known in literature,
%
\begin{equation}
    k_{\mathrm{RBF}}(x, x')
    = \sigma_f^2 \exp\!\left(
        -\frac{\lVert x - x' \rVert^2}{2\ell^2}
      \right) .
    \label{eq:rbf-kernel}
\end{equation}
%
Behavior of this kernel is controlled by two hyperparameters - the signal variance $\sigma_f^2$ sets the vertical amplitude of the sampled functions, while the length-scale $\ell$ sets the distance over whichfunction values remain correlated: a small $\ell$ produces rapidly varying functions, whereas a large $\ell$ produces nearly flat ones. Because the correlation decays with distance, the prior is recovered far away from the training inputs. This decay is the mechanism by which a GP increases its epistemic uncertainty under extrapolation, and it is examined experimentally in Chapter~5.

Observation noise is included by adding a white-noise term to the kernel,
%
\begin{equation}
    k_y(x, x') = k_{\mathrm{RBF}}(x, x') + \sigma_n^2\, \delta_{xx'},
    \label{eq:noisy-kernel}
\end{equation}
%
where $\delta_{xx'}$ equals one for identical set of inputs and zero otherwise. The choice of kernel is a strong modelling assumption rather than a technical detail. Samples from an RBF prior are infinitely differentiable, which makes the kernel a poor description of rough functions; the
Mat\'ern class relaxes this requirement at the cost of an additional hyperparameter~\cite{rasmussen2006}. The RBF kernel with a single shared length-scale was
adopted here due to the fact it is the standard baseline in the literature on uncertainty estimation, and tuning a separate kernel per dataset would make the comparison between methods
dependent on the effort spent on the reference model.

% [Miejsce na Rysunek 3.1: wpływ length-scale na prior.
%  Trzy panele obok siebie, w każdym 5 próbek funkcji z GP prior na przedziale
%  [-2, 10] dla ell = 0.3, 1.0, 3.0 (sigma_f = 1). Cel rysunku: pokazać, że
%  kernel koduje założenia o gładkości, zanim pojawią się jakiekolwiek dane.]

\begin{figure}[htbp]
    \centering
    \includegraphics[width=0.95\textwidth]{rodzial3_rys/img3_1.png}
    \caption{Samples drawn from a Gaussian process prior with an RBF kernel for three values
    of the length-scale $\ell$ (signal variance fixed at $\sigma_f^2 = 1$). Short length-scales
    admit rapidly varying functions, while long ones admit only slowly varying ones. The kernel
    encodes assumptions about smoothness before any data are observed.}
    \label{fig:gp-prior-samples}
\end{figure}

% ---------------------------------------------------------------------
\subsection{Posterior and predictive distribution}
\label{subsec:gp-posterior}

Let $\mathbf{y}$ denote the vector of observed targets and $x^\ast$ a test input. Under the
generative model of Equation~\eqref{eq:generative-model} with Gaussian noise, the observed
targets and the function value at the test point are jointly Gaussian,
%
\begin{equation}
    \begin{bmatrix} \mathbf{y} \\ f(x^\ast) \end{bmatrix}
    \sim
    \mathcal{N}\!\left(
      \mathbf{0},
      \begin{bmatrix}
        K + \sigma_n^2 I & \mathbf{k}_\ast \\
        \mathbf{k}_\ast^{\!\top} & k(x^\ast, x^\ast)
      \end{bmatrix}
    \right),
    \qquad
    \left[\mathbf{k}_\ast\right]_i = k(x_i, x^\ast) .
    \label{eq:gp-joint}
\end{equation}
%
Conditioning the joint distribution on the observations yields the predictive distribution in
closed form,
%
\begin{equation}
    \mu_\ast
    = \mathbf{k}_\ast^{\!\top}
      \left(K + \sigma_n^2 I\right)^{-1} \mathbf{y},
    \label{eq:gp-posterior-mean}
\end{equation}
%
\begin{equation}
    \sigma_\ast^2
    = k(x^\ast, x^\ast)
      - \mathbf{k}_\ast^{\!\top}
        \left(K + \sigma_n^2 I\right)^{-1} \mathbf{k}_\ast .
    \label{eq:gp-posterior-var}
\end{equation}
%
Equations~\eqref{eq:gp-posterior-mean} and~\eqref{eq:gp-posterior-var} are the exact
counterpart of the intractable integral~\eqref{eq:predictive-2}.  The predictive mean is a linear combination of the observed targets, weighted by the similarity between the test point and each training point. The predictive variance, does not depend on $\mathbf{y}$ at all. A GP therefore derives its epistemic uncertainty from the geometry of the training inputs, before any target value is taken into account. The second term of Equation~\eqref{eq:gp-posterior-var} is the reduction in prior variance contributed by the data, and it fades away where no data are available.

The predictive variance of a noisy observation is
%
\begin{equation}
    \operatorname{Var}\!\left(y^\ast \mid x^\ast\right)
    = \underbrace{\sigma_\ast^2}_{\text{epistemic}}
      + \underbrace{\sigma_n^2}_{\text{aleatoric}} ,
    \label{eq:gp-variance-split}
\end{equation}
%
so both components are read directly from the model rather than estimated by sampling. 
The white-noise term of Equation~\eqref{eq:noisy-kernel} is constant across the input space, which
makes the model homoscedastic, in agreement with the noise model adopted by every other method
compared in this work.

% [Miejsce na Rysunek 3.2: posterior GP na zbiorze load_sin.
%  Trening na [0, 6], ewaluacja na [-2, 10]. Predykcyjna średnia, pasmo +-2 sigma,
%  punkty treningowe, prawdziwa funkcja linią przerywaną. Cel rysunku: wąskie pasmo
%  w zakresie treningowym, rozejście się poza nim, powrót do prior mean = 0.]

\begin{figure}[htbp]
    \centering
    \includegraphics[width=0.85\textwidth]{rodzial3_rys/img3_2.png}
    \caption{Gaussian process posterior on the synthetic sine dataset. The model is trained on
    the interval $[0, 6]$ and evaluated on $[-2, 10]$. The $\pm 2\sigma_\ast$ band is narrow
    where training points are dense and widens rapidly outside the training range, where the
    predictive mean reverts towards the prior mean.}
    \label{fig:gp-posterior-sine}
\end{figure}

% ---------------------------------------------------------------------
\subsection{Hyperparameter selection}
\label{subsec:gp-hyperparameters}

The hyperparameters $\vartheta = (\sigma_f, \ell, \sigma_n)$ are not known in advance but
estimated from the data. The criterion is the log marginal likelihood,
%
\begin{equation}
    \log p(\mathbf{y} \mid X, \vartheta)
    = -\tfrac{1}{2}\, \mathbf{y}^{\!\top} K_y^{-1} \mathbf{y}
      - \tfrac{1}{2}\, \log \lvert K_y \rvert
      - \tfrac{N}{2} \log 2\pi ,
    \qquad
    K_y = K + \sigma_n^2 I .
    \label{eq:gp-marginal-likelihood}
\end{equation}
%
The first term rewards the fit to the data, while the second penalises model complexity through the
determinant of the covariance matrix, and the third one is a constant. The quantity being
maximised is the evidence $p(\mathcal{D})$ from Equation~\eqref{eq:bayes-rule}, which was
identified in Section~\ref{subsec:intractability} as the source of every computational
difficulty. For a GP it is available analytically, and the trade-off between fit and
complexity is therefore resolved by the model itself, without a validation set.

The objective in Equation~\eqref{eq:gp-marginal-likelihood} is not convex and may possess
several local optima. The optimiser is consequently restarted from several random
initialisations, and the solution with the highest marginal likelihood is retained. It should
be noted that this procedure returns point estimates of the hyperparameters, an approach known
as type~II maximum likelihood. A fully Bayesian treatment would also marginalize over $\vartheta$. The GP used here is thus exact given the kernel, and one small component of uncertainty,
namely uncertainty about the kernel itself, remains outside the reported  intervals~\cite{rasmussen2006}.

% ---------------------------------------------------------------------
\subsection{Computational cost and practical limitations}
\label{subsec:gp-limitations}

Both the predictive equations and the marginal likelihood require the inversion of an
$N \times N$ matrix, usually performed through a Cholesky decomposition. The cost is
$\mathcal{O}(N^3)$ in time and $\mathcal{O}(N^2)$ in space, and it is incurred anew at every
step of hyperparameter optimisation. Exact GP regression is therefore restricted to datasets of
moderate size. For this reason, the benchmark training sets in this work were subsampled to a
fixed size, so that every method is trained on identical data and the comparison is not
confounded by the scaling limits of a single model. Sparse approximations based on inducing
points reduce the cost substantially~\cite{quinonero2005}, but they were excluded, because they
would replace the exact posterior with an approximate one.

Two further limitations should be noted: stationary kernels such as the RBF depend on the
distance between inputs and distances become uninformative as the input dimension grows, which
weakens the method on high-dimensional problems. In addition, the reversion of the predictive
mean to the prior mean far from the data is desirable for uncertainty estimation, but it
provides no extrapolation of the trend itself.


%\section{Bayesian neural networks (Bayes by Backprop)}
\section{Bayesian neural networks (Bayes by Backprop)}
\label{sec:bayes-by-backprop}

Section~\ref{subsec:variational-inference} introduced variational inference as a
general strategy rather than as a particular algorithm. Bayes by Backprop (BBB) is
one concrete instance of that strategy, applied to the weights of a neural
network~\cite{blundell2015}. The contrast with Section~\ref{sec:gaussian-processes}
is based on the assumption that the Gaussian process performs inference in the function space and obtains the predictive distribution in closed form, whereas the BBB performs inference in
the parameter space and settles for an approximation. The following therefore describes a
realization of the mathematical foundations of Chapter~\ref{ch:methods}, and the derivations of
Section~\ref{subsec:variational-inference} are not repeated.

Each weight of the network is replaced by a pair of variational parameters instead
of a single number $w_j$, the model stores a mean $\mu_j$ and a scale parameter
$\rho_j$, from which the standard deviation is obtained as   $\sigma_j = \log\!\left(1 + \exp \rho_j\right)$. The softplus transform keeps $\sigma_j$ positive without a constraint on the optimizer. The resulting variational family is exactly the mean-field Gaussian of Equation~\eqref{eq:mean-field}, with $\theta = (\mu, \rho)$ collecting all variational parameters. Under this parameterization, training a single network is replaced by training an infinite
ensemble whose members draw their weights from one shared, learnt distribution~\cite{blundell2015}.

% [Miejsce na Rysunek 3.4: waga deterministyczna kontra waga jako rozkład.
%  Dwa panele obok siebie, ten sam mały graf sieci (2-3-1). Po lewej każda
%  krawędź opisana jedną liczbą, po prawej — małą krzywą gaussowską o środku
%  mu_j i szerokości sigma_j. Cel rysunku: pokazać bez wzoru, na czym polega
%  różnica między punktowym oszacowaniem a rozkładem nad wagami.
%  UWAGA: rysunek narysować SAMODZIELNIE. Odpowiedniki u Blundella (rys. 1)
%  i Jospina (rys. 3) są objęte prawem autorskim i nie wolno ich przedrukować.]

\begin{figure}[htbp]
    \centering
    \includegraphics[width=0.85\textwidth]{rodzial3_rys/img3_4.png}
    \caption{Weight representation in a deterministic network and in a Bayesian
    neural network trained with Bayes by Backprop. On the left, each connection
    carries a single value obtained by point estimation. On the right, each
    connection carries a Gaussian distribution described by a mean $\mu_j$ and a
    standard deviation $\sigma_j$, so the number of stored parameters doubles.}
    \label{fig:bbb-weight-distributions}
\end{figure}

The objective function follows directly from Equation~\eqref{eq:elbo}, rewritten in
the notation used for network weights. Maximizing the evidence lower bound is
equivalent to minimizing the variational free energy,
%.
\begin{equation}
    \mathcal{F}(\mathcal{D}, \theta)
    = \underbrace{\mathrm{KL}\bigl(q(\mathbf{w} \mid \theta) \,\|\, p(\mathbf{w})\bigr)}_{\text{complexity cost}}
      - \underbrace{\mathbb{E}_{q(\mathbf{w} \mid \theta)}
        \bigl[\log p(\mathcal{D} \mid \mathbf{w})\bigr]}_{\text{likelihood cost}} ,
    \label{eq:bbb-free-energy}
\end{equation}
%
where $\mathbf{w}$ denotes a sample of all network weights~\cite{blundell2015}. The
likelihood cost rewards are in agreement with the training data. The complexity cost
penalizes departure from the prior and admits an interpretation as a compression
cost in the sense of minimum description length.

The second term of Equation (complexity cost) ~\eqref{eq:bbb-free-energy} deserves emphasis, because it settles a question raised in Section~\ref{subsec:mle-map}. Weight decay was shown
there to be MAP estimation under a Gaussian prior, that is, a regularization term
constructed on fact from a Bayesian reading. Here, the relationship is reversed - the penalty is not added to the loss by design but follows from the variational objective itself. Regularisation in a Bayesian neural network is therefore derived rather than imposed, and the strength of the penalty is fixed by the choice of prior rather than by a tuned coefficient~\cite{blundell2015}.

% ---------------------------------------------------------------------
\subsection{Training with unbiased Monte Carlo gradients}
\label{subsec:bbb-training}

The obstacle to optimizing Equation~\eqref{eq:bbb-free-energy} is not the objective, but the gradient. The weight sampling introduces stochastic nodes inside the network, and the backpropagation cannot pass through them in the usual way~\cite{jospin2022}. The expectation in the likelihood cost is taken over a distribution that itself depends on the parameters being optimised, so the gradient  cannot be moved inside the expectation without further argument. The
reparameterisation approach of Equation~\eqref{eq:reparameterization} removes the difficulty by writing a weight sample as a deterministic function of the variational parameters and of parameter-free noise~\cite{kingma2014}. 
For the Gaussian family adopted here, this takes the form
%
\begin{equation}
    \mathbf{w} = \mu + \log\!\left(1 + \exp \rho\right) \odot \varepsilon ,
    \qquad \varepsilon \sim \mathcal{N}(0, I) ,
    \label{eq:bbb-reparameterisation}
\end{equation}
%
where $\odot$ denotes multiplication in the element (elementwise). Randomness is limited to
$\varepsilon$, and gradients with respect to $\mu$ and $\rho$ pass through ordinary
backpropagation.

Variational treatments of neural networks predate this construction. Blundell et
al.\ build on the work of Graves (2011) and show that the gradients used there can
be made unbiased, and that the method can be extended to non-Gaussian priors~\cite{blundell2015}. Both improvements evolved from a single decision: the complexity cost is not evaluated in closed form. The objective is instead estimated entirely by sampling,
%
\begin{equation}
    \mathcal{F}(\mathcal{D}, \theta)
    \approx \sum_{i=1}^{n}
      \left[
        \log q\bigl(\mathbf{w}^{(i)} \mid \theta\bigr)
        - \log p\bigl(\mathbf{w}^{(i)}\bigr)
        - \log p\bigl(\mathcal{D} \mid \mathbf{w}^{(i)}\bigr)
      \right] ,
    \label{eq:bbb-mc-cost}
\end{equation}
%
where $\mathbf{w}^{(i)}$ is the $i$-th sample drawn from the variational posterior.
Dispensing with the closed form frees the choice of prior. Blundell et al.\ exploit
this freedom with a scale mixture of two zero-mean Gaussians, the first broad and
the second sharply concentrated around zero. On a network with two hidden layers of
1200 units, this prior reached a test error of 1.32\% on MNIST, against 2.04\% for a
plain Gaussian prior.

Equation~\eqref{eq:bbb-mc-cost} has another less obvious property - every term is
evaluated on the same weight sample, a variance reduction technique known as common
random numbers~\cite{blundell2015}. Where part of the objective is available in
closed form, that part is insensitive to the particular draw, while the remainder is
not. A qualification is in order, and it comes from the authors themselves. The
sampled complexity cost was not found to outperform a closed-form Kullback--Leibler
term where the latter could be computed, only to perform no worse. The argument for
the sampling estimator rests on the range of admissible priors, not on measured
accuracy.

The resulting training loop remains close to that of a deterministic
network~\cite{jospin2022}. Each step draws $\varepsilon$, assembles the weights
through Equation~\eqref{eq:bbb-reparameterisation}, evaluates the objective of
Equation~\eqref{eq:bbb-mc-cost} on a minibatch, and updates $\mu$ and $\rho$ by
gradient descent. The gradient with respect to the sampled weights is shared between
both updates and is precisely what ordinary backpropagation returns~\cite{blundell2015}.
Standard optimizers such as Adam apply without modification. One further detail
concerns minibatches: the complexity cost refers to the whole data set and must be
distributed between them, for instance, by weighting it with coefficients
$\pi_i \in [0,1]$ satisfying $\sum_{i=1}^{M} \pi_i = 1$,
%
\begin{equation}
    \mathcal{F}_i^{\pi}(\mathcal{D}_i, \theta)
    = \pi_i\, \mathrm{KL}\bigl(q(\mathbf{w} \mid \theta) \,\|\, p(\mathbf{w})\bigr)
      - \mathbb{E}_{q(\mathbf{w} \mid \theta)}
        \bigl[\log p(\mathcal{D}_i \mid \mathbf{w})\bigr] .
    \label{eq:bbb-kl-reweighting}
\end{equation}
%
A schedule that front-loads the complexity cost onto the first minibatches of each
epoch was reported to work well~\cite{blundell2015}. It should be noted that the
objective is estimated from a small number of samples, often a single one, so the
gradient estimate is noisy and the convergence curve is far less smooth than in
deterministic training~\cite{jospin2022}. Averaging the reported loss over several
epochs gives a more readable approach of convergence.

% ---------------------------------------------------------------------
\subsection{Prediction and uncertainty decomposition}
\label{subsec:bbb-prediction}

The prediction reproduces the marginalization of Equation~\eqref{eq:predictive-2} by
Monte Carlo integration. A test input is passed through the network $T$ times, and
each pass draws an independent weight sample $\mathbf{w}^{(t)}$ from the variational
posterior. 
Each pass yields a predictive mean $\mu_{\mathbf{w}^{(t)}}(x^\ast)$ and a
predictive variance $\sigma^2_{\mathbf{w}^{(t)}}$, which under the homoscedastic
noise model adopted here does not depend on the input. The prediction reported for $x^\ast$ is the
average across passes~\cite{jospin2022},
%
\begin{equation}
    \hat{\mu}(x^\ast)
    = \frac{1}{T} \sum_{t=1}^{T} \mu_{\mathbf{w}^{(t)}}(x^\ast) .
    \label{eq:bbb-predictive-mean}
\end{equation}

The spread across passes is not a by-product of this averaging but the quantity of
primary interest. Both components of Equation~\eqref{eq:total-variance-2} are
estimated from the same $T$ samples,
%
\begin{equation}
    \operatorname{Var}\!\left(y^\ast \mid x^\ast\right)
    \approx
    \underbrace{\frac{1}{T} \sum_{t=1}^{T}
      \sigma^2_{\mathbf{w}^{(t)}}(x^\ast)}_{\text{aleatoric}}
    +
    \underbrace{\frac{1}{T-1} \sum_{t=1}^{T}
      \bigl(\mu_{\mathbf{w}^{(t)}}(x^\ast) - \hat{\mu}(x^\ast)\bigr)^2}_{\text{epistemic}} .
    \label{eq:bbb-variance-estimator}
\end{equation}
%
The first term averages the noise predicted by the individual networks, while the second
measures how far those networks disagree about the mean, and it vanishes when the
variational posterior collapses to a point. Where the training data are dense, the
sampled networks converges and the epistemic term is small, but if there is absent data, the
weight distributions are constrained only by the prior, the sampled networks diverge, and the epistemic term grows.

% [Miejsce na Rysunek 3.5: posterior BBB na zbiorze load_sin.
%  Trening na [0, 6], ewaluacja na [-2, 10]. TE SAME osie i ta sama skala co
%  Rysunek 3.2 (posterior GP), żeby oba dały się porównać wzrokiem.
%  Predykcyjna średnia z Równania (bbb-predictive-mean), pasmo +-2 sigma
%  z pierwiastka Równania (bbb-variance-estimator), punkty treningowe,
%  prawdziwa funkcja linią przerywaną. Opcjonalnie: 10-15 pojedynczych
%  przebiegów cienką linią, żeby widać było rozrzut, z którego liczony jest
%  człon epistemiczny. Cel rysunku: czy pasmo rozchodzi się poza zakresem
%  treningowym słabiej niż u GP.]

\begin{figure}[htbp]
    \centering
    \includegraphics[width=0.85\textwidth]{rodzial3_rys/img3_5.png}
    \caption{Predictive distribution of a Bayesian neural network trained with Bayes
    by Backprop on the synthetic sine dataset. The model is trained on the interval
    $[0, 6]$ and evaluated on $[-2, 10]$, with the axes matching those of
    Figure~\ref{fig:gp-posterior-sine} to allow a direct comparison. The band shows
    $\pm 2\sigma$ obtained from Equation~\eqref{eq:bbb-variance-estimator} over
    $T$ stochastic forward passes.}
    \label{fig:bbb-posterior-sine}
\end{figure}

% ---------------------------------------------------------------------
\subsection{Computational cost and limitations}
\label{subsec:bbb-limitations}
 
The cost of the method is incurred twice. Training stores two variational parameters
per weight, so the parameter count doubles relative to a deterministic network of the
same architecture~\cite{arbel2023}. Each training iteration was reported to run about
twice as slowly as an iteration with dropout~\cite{blundell2015}. Inference is
affected more severely because a single forward pass no longer suffices and the
$T$ passes of Equation~\eqref{eq:bbb-variance-estimator} must be performed for every
test point. The need for Monte Carlo integration at evaluation time is a persistent
drawback of Bayesian neural networks and not a property of this algorithm
alone~\cite{jospin2022}.

The mean-field assumption is the most consequential limitation. Learning a complete
covariance matrix on the weights would require $\mathcal{O}(n^2)$ variational
parameters, which is unattainable for the parameter counts reported in
Section~\ref{subsec:intractability}; the diagonal approximation reduces this to
$\mathcal{O}(n)$~\cite{jospin2022}. The independence between weights is therefore a
concession to tractability rather than a modeling belief. Its cost is stated
plainly in the literature: the restriction is likely to result in an underestimate
of the overall posterior uncertainty~\cite{arbel2023}. Variational methods are also
prone to mode collapse, that is, to concentrating on a single mode of a posterior
known to be highly multimodal. Both effects act in the same direction and narrow the
approximate posterior.

A qualification is again due, and again it comes from the original work. Blundell et
al.\ acknowledge the known tendency of variational methods to underestimate
uncertainty, yet report that they did not observe its practical consequences in
their experiments~\cite{blundell2015}. The claim is therefore a prediction rather
than a settled result. Equation~\eqref{eq:total-variance-2} identifies where such an
underestimate would appear: the epistemic term of Equation~\eqref{eq:bbb-variance-estimator} is the only component affected by the shape of the approximate posterior. This gives a testable expectation for Chapter~\ref{ch:results}. If the mean-field posterior is too narrow, Bayes by
Backprop should produce narrower prediction intervals than the Gaussian process of
Section~\ref{sec:gaussian-processes}, and its empirical coverage should fall short of
the nominal level. The metrics defined in Chapter~\ref{ch:setup} are reported for
both methods precisely to settle this question.


%\section{Monte carlo dropout}
\section{Monte Carlo dropout}
\label{sec:mc-dropout}

Section~\ref{sec:bayes-by-backprop} obtained an approximate posterior at the cost
of doubling the number of stored parameters, while the Monte Carlo (MC) dropout (presented below) obtains one at no cost at all. Gal and Ghahramani showed that a network trained with dropout
before every weight layer, together with $L_2$ regularisation, is already performing approximate variational inference; nothing in the training procedure needs to change, and the only modification occurs at test time~\cite{gal2016}. The method therefore belongs to the same family as Bayes by Backprop, and differs from it in a fact that the variational family is
Bernoulli rather than Gaussian.

The starting point is the objective minimized by an ordinary dropout network. For
a model with $L$ weight layers, weight matrices $\mathbf{W}_i$, bias vectors
$\mathbf{b}_i$ and a loss $E(\cdot,\cdot)$ such as the squared error, weight decay
$\lambda$ gives
%
\begin{equation}
    \mathcal{L}_{\mathrm{dropout}}
    = \frac{1}{N} \sum_{n=1}^{N} E\bigl(y_n, \hat{y}_n\bigr)
      + \lambda \sum_{i=1}^{L}
        \left( \lVert \mathbf{W}_i \rVert_2^2
             + \lVert \mathbf{b}_i \rVert_2^2 \right) .
    \label{eq:dropout-objective}
\end{equation}
%
Dropout itself is applied by sampling a binary vector for every unit of every
layer and zeroing the corresponding rows of the weight matrix. Written as a
distribution over the weights, this defines the variational family
%
\begin{equation}
    \mathbf{W}_i = \mathbf{M}_i \operatorname{diag}
      \bigl( [z_{i,j}]_{j=1}^{K_i} \bigr),
    \qquad
    z_{i,j} \sim \mathrm{Bernoulli}(p_i) ,
    \label{eq:bernoulli-family}
\end{equation}
%
in which the matrices $\mathbf{M}_i$ are the variational parameters and $p_i$ is
the probability that a unit is retained~\cite{gal2016}. The contrast with the
mean-field Gaussian family of Equation~\eqref{eq:mean-field} is sharp: each weight
of Equation~\eqref{eq:bernoulli-family} takes one of two values, either $m_{i,jk}$
or zero, so the approximate posterior is a discrete distribution rather than a
continuous one.

The central result of~\cite{gal2016} is that minimizing Equation~\eqref{eq:dropout-objective} is equivalent, up to a positive scaling constant, to minimizing the Kullback--Leibler divergence of
Equation~\eqref{eq:vi-objective} between the family \eqref{eq:bernoulli-family} and the posterior of a deep Gaussian process marginalized over its covariance function parameters resulting in two steps of the derivation account for the two terms of Equation~\eqref{eq:dropout-objective}. The likelihood term of the evidence lower bound \eqref{eq:elbo} is estimated by Monte Carlo integration with a single sample per data point, which is precisely what a dropout mask provides per forward pass. The divergence from the prior admits, for a large number of hidden units, a closed-form approximation proportional to $\lVert \mathbf{M}_i \rVert_2^2$, recovering the weight-decay  term~\cite{gal2016appendix}.

This result closes a loop opened in Section~\ref{sec:gaussian-processes} -the
Gaussian process was introduced there as the reference method, the only one whose
predictive distribution is exact; here it reappears as the model being approximated, this time in its deep variant. The practical consequence is stated clearly by the authors: epistemic information is already contained in a trained dropout network and has been discarded so far~\cite{gal2016} by the standard test-time procedure, and recovery of the information
requires only that this procedure be replaced.

% ---------------------------------------------------------------------
\subsection{Prediction and uncertainty decomposition}
\label{subsec:mc-dropout-prediction}

At test time, an ordinary dropout network scales each weight matrix by its retention probability $p_i$ and performs one deterministic forward pass. MC dropout instead leaves the sampling switched on and performs $T$ stochastic passes, each with an independently drawn set of masks $\boldsymbol{\omega}^{(t)}$. The predictive mean is their average,
%
\begin{equation}
    \hat{\mu}(x^\ast)
    = \frac{1}{T} \sum_{t=1}^{T}
      \hat{y}\bigl(x^\ast, \boldsymbol{\omega}^{(t)}\bigr) ,
    \label{eq:mc-dropout-mean}
\end{equation}
%
which is a Monte Carlo estimate of the marginalization in Equation~\eqref{eq:predictive-2}. It should be noted that Equation~\eqref{eq:mc-dropout-mean} is not itself a new proposal. Averaging
several stochastic passes was already known as model averaging, and the  weight-scaling rule was justified empirically as an approximation to it. What Bayesian reading adds is a derivation, and with it access to the higher moments of the predictive distribution~\cite{gal2016}.

The predictive variance follows by moment matching over the same $T$ samples,
%
\begin{equation}
    \operatorname{Var}\!\left(y^\ast \mid x^\ast\right)
    \approx
    \underbrace{\tau^{-1}}_{\text{aleatoric}}
    +
    \underbrace{\frac{1}{T-1} \sum_{t=1}^{T}
      \bigl( \hat{y}(x^\ast, \boldsymbol{\omega}^{(t)})
             - \hat{\mu}(x^\ast) \bigr)^2}_{\text{epistemic}} ,
    \label{eq:mc-dropout-variance}
\end{equation}
%
where $\tau$ is the model precision, which is not a free parameter, but is
fixed by the quantities already present in Equation~\eqref{eq:dropout-objective}, through the identity
%
\begin{equation}
    \tau = \frac{p\, l^2}{2 N \lambda} ,
    \label{eq:mc-dropout-precision}
\end{equation}
%
in which $l$ is a prior length-scale expressing a belief about the frequency of the modeled function~\cite{gal2016appendix}. Equations \eqref{eq:mc-dropout-variance} and \eqref{eq:mc-dropout-precision} expose a limitation of the method in its original form. The aleatoric term is a single
constant across the whole input space (what makes MC dropout a homoscedastic model) and the value of that constant is inherited from the dropout rate, the weight decay and a user-specified length-scale rather than estimated from the data. The situation mirrors the white-noise term of Equation~\eqref{eq:noisy-kernel} and the decomposition \eqref{eq:gp-variance-split}, with one difference: the former is fitted by maximising the marginal likelihood, whereas $\tau^{-1}$ is
not fitted at all.

A later modification of the method removes this limitation by combining dropout
sampling with the two-headed network of
Section~\ref{subsec:aleatoric}~\cite{kendall2017}.

Each stochastic pass then returns a mean $\hat{\mu}_{\boldsymbol{\omega}^{(t)}}(x^\ast)$ and a variance
$\hat{\sigma}^2_{\boldsymbol{\omega}^{(t)}}(x^\ast)$, the network is trained with the Gaussian negative logarithmic likelihood of Equation~\eqref{eq:gaussian-nll}, and the constant $\tau^{-1}$ of Equation~\eqref{eq:mc-dropout-variance} is replaced by an average of predicted noise variances,
%
\begin{equation}
    \operatorname{Var}\!\left(y^\ast \mid x^\ast\right)
    \approx
    \underbrace{\frac{1}{T} \sum_{t=1}^{T}
      \hat{\sigma}^2_{\boldsymbol{\omega}^{(t)}}(x^\ast)}_{\text{aleatoric}}
    +
    \underbrace{\frac{1}{T-1} \sum_{t=1}^{T}
      \bigl( \hat{\mu}_{\boldsymbol{\omega}^{(t)}}(x^\ast)
             - \hat{\mu}(x^\ast) \bigr)^2}_{\text{epistemic}} .
    \label{eq:mc-dropout-heteroscedastic}
\end{equation}
%
Equation~\eqref{eq:mc-dropout-heteroscedastic} is the empirical counterpart of
Equation~\eqref{eq:total-variance-2} and is identical in form to the estimator
\eqref{eq:bbb-variance-estimator} of Section~\ref{sec:bayes-by-backprop}. The
aleatoric term is no longer inherited from the regularisation hyperparameters
but read from a variance head trained by the objective of
Equation~\eqref{eq:gaussian-nll}. The two variants of
Equations~\eqref{eq:mc-dropout-variance}
and~\eqref{eq:mc-dropout-heteroscedastic} therefore differ in how the aleatoric
term is obtained, and the choice between them is a property of the noise model
rather than of the posterior approximation.

% Uwaga dla autora: Kendall i Gal zapisują człon epistemiczny w postaci
% (1/T) sum y_t^2 - mu^2, czyli z dzielnikiem T. Powyżej użyto dzielnika T-1,
% spójnie z Równaniem (bbb-variance-estimator) z Sekcji 3.2. Różnica jest
% pomijalna dla używanych T, ale wybór trzeba zadeklarować raz w Rozdziale 4
% i stosować dla wszystkich metod.

The two terms behave differently as the test point moves away from the training
data, and this behavior has been measured. Authors report that the mean aleatoric variance remains approximately constant when the training set is reduced or when the model is evaluated on an unrelated dataset, whereas the epistemic variance falls as the training set grows and rises sharply outside the distribution~\cite{kendall2017}. This is the empirical form of the distinction
drawn in Section~\ref{subsec:epistemic}, and it is the property that Equation~\eqref{eq:mc-dropout-heteroscedastic} is meant to expose.

%% ryas 3.6 do zmiany raczej - lekka zmiana teorii !!!!!!!!!!!!!!!!!!!!!!!!!
% [Miejsce na Rysunek 3.6: posterior MC dropout na zbiorze load_sin.
%  Trening na [0, 6], ewaluacja na [-2, 10]. TE SAME osie i ta sama skala co
%  Rysunek 3.2 (posterior GP) i Rysunek 3.5 (posterior BBB) — trzy rysunki
%  mają dać się porównać wzrokiem, to jest główny argument wizualny rozdziału.
%  Predykcyjna średnia z Równania (mc-dropout-mean), pasmo +-2 sigma
%  z pierwiastka Równania (mc-dropout-heteroscedastic), punkty treningowe,
%  prawdziwa funkcja linią przerywaną. Opcjonalnie: 10-15 pojedynczych
%  przebiegów cienką linią. Sieć: MLP z src/models.py z dropout > 0,
%  z model.train() przy predykcji, T = 100 przebiegów.
%  Cel rysunku: czy pasmo rozchodzi się poza zakresem treningowym, i o ile
%  słabiej niż u GP.]

\begin{figure}[htbp] 
    \centering
    \includegraphics[width=0.85\textwidth]{rodzial3_rys/img3_6.png}
    \caption{Predictive distribution of a network evaluated with Monte Carlo
    dropout on the synthetic sine dataset. The model is trained on the interval
    $[0, 6]$ and evaluated on $[-2, 10]$, with the axes matching those of
    Figures~\ref{fig:gp-posterior-sine} and~\ref{fig:bbb-posterior-sine} to allow
    a direct comparison. The band shows $\pm 2\sigma$ obtained from
    Equation~\eqref{eq:mc-dropout-heteroscedastic} over $T$ stochastic forward
    passes.}
    \label{fig:mc-dropout-posterior-sine}
\end{figure}

% ---------------------------------------------------------------------
\subsection{Computational cost and limitations}
\label{subsec:mc-dropout-limitations}

The main adventage of the method is the fact that it does not require any additional parameters added to the network, so the storage doubling described in Section~\ref{subsec:bbb-limitations} does not occur, and an existing model trained with dropout can be re-evaluated without retraining~\cite{gal2016}. The cost appears in inference alone, where $T$ passes replace one. Gal and Ghahramani observe that the passes are independent and may be executed concurrently, which
would leave the running time unchanged. Here, a qualification is due, and it comes from a later application of the method. Kendall and Gal report a fiftyfold slow-down for fifty samples on a dense architecture, and attribute it to memory constraints that prevent the passes from being parallelised~\cite{kendall2017}. The method is therefore free in parameters but not in evaluation time.

The dropout rate is the more delicate issue - in Equation~\eqref{eq:bernoulli-family} the probability $p_i$ is a property of the variational family, and yet it is set by hand and held fixed throughout training. The spread of the approximate posterior is thus chosen rather than inferred. Two boundary cases make the point: at $p = 1$ the family collapses to a point mass and the method reduces to the MAP estimation of Section~\ref{subsec:mle-map}, while at $p = 0$ all the mass is placed at the origin~\cite{lefolgoc2021}. Gal and Ghahramani report that models initialized with rates of $0.1$ and $0.2$ differ in their uncertainty early in training, but become almost indistinguishable once converged~\cite{gal2016}. However, the rate is a hyperparameter that governs an uncertainty estimate, and it is treated as such in Chapter~\ref{ch:setup}.

A second discrepancy separates the theory from common practice, the derivation requires dropout before every weight layer, whereas implementations frequently apply it to the final layers alone. This network interleaves MAP estimates in the layers without dropout with approximately Bayesian ones in the layers with it~\cite{gal2016appendix}. The uncertainty it reports is, consequently, not the uncertainty the derivation describes.

The quality of the estimates is bounded by the same argument that justifies the method. Because a dropout network approximates a Gaussian process, it inherits the properties of that process, including the fact that its uncertainty is not calibrated~\cite{gal2016appendix}. Section~\ref{subsec:gp-kernel} noted that the choice of covariance function is a modeling assumption; under the present correspondence, that assumption is made by choosing a non-linearity. The effect is visible in extrapolation. Gal and Ghahramani observe that the uncertainty continues to grow away from the data for a ReLU network, whereas it remains bounded for a TanH network, because the two saturate differently and correspond to different covariance functions~\cite{gal2016}. The activation function is therefore not a neutral implementation detail once the uncertainty is the quantity of interest, which is why it is kept fixed across all methods in Chapter~\ref{ch:setup}.

The strongest objection concerns the Bayesian status of the method itself. Le Folgoc et al.\ observe that the family \eqref{eq:bernoulli-family} places all its mass on a finite set of $2^P$ parameter configurations, obtained from a single vector $\hat{\theta}$ by zeroing subsets of its coordinates~\cite{lefolgoc2021}. Three consequences follow:
\begin{itemize}
    \item On a regression problem with a posterior available in closed form, the true parameter vector receives zero probability under the  approximation, as it does under MAP estimation and unlike any Gaussian family.
    \item The multimodality of the approximate posterior, sometimes presented as an advantage over mean-field methods, is an artifact of this construction rather than a property of the target: the modes sit at the projections of $\hat{\theta}$ onto the coordinate hyperplanes, and their location carries no information about the data.
    \item Since the whole distribution has one $P$-dimensional degree of freedom, it cannot represent a genuinely multimodal posterior
\end{itemize}
The authors propose reading MC dropout as a modification of the Bayesian model, by a sparsity-inducing prior, rather than as an approximation to it~\cite{lefolgoc2021}.

This objection is not settled in the literature and the present work does not attempt to settle it. However, it does determine how the results are read. The quantity reported for MC dropout in Chapter~\ref{ch:results} is the empirical spread of $T$ stochastic passes, which remains well defined whichever interpretation is adopted; the claim that this spread approximates a posterior
variance is the contested part . A concrete expectation can be attached to it: variational methods under-estimate posterior spread unless the true posterior belongs to the assumed family, and Gal and Ghahramani themselves report that on an interpolation task MC dropout was over-confident, while the Gaussian process over-estimated its interval~\cite{gal2016appendix}. The prediction stated for Bayes by Backprop in Section~\ref{subsec:bbb-limitations} therefore extends to MC
dropout, and for the same reason: prediction intervals narrower than those of the Gaussian process and empirical coverage below the nominal level. The PICP and MPIW metrics of Chapter~\ref{ch:setup} are reported for both methods to test it.

% [Miejsce na Rysunek 3.7: wpływ liczby przebiegów T na oszacowanie niepewności.
%  Zbiór load_sin, ta sama sieć co na Rysunku 3.6. Wariant A (preferowany):
%  średnia szerokość pasma +-2 sigma oraz RMSE w funkcji T = 2, 5, 10, 20, 50,
%  100, 200, 500, każdy punkt uśredniony po 10 losowaniach masek, z pasmem
%  odchylenia standardowego między losowaniami. Cel rysunku: pokazać, od którego
%  T oszacowanie się stabilizuje, i uzasadnić liczbowo wartość T przyjętą
%  w Rozdziale 4 — bez tego wybór T wygląda na arbitralny.
%  Wariant B (jeśli A wyjdzie nieczytelnie): trzy panele z pasmem +-2 sigma dla
%  T = 5, 25, 100 na tych samych osiach co Rysunek 3.6.]

\begin{figure}[htbp]
    \centering
    \includegraphics[width=0.85\textwidth]{rodzial3_rys/img3_7.png}
    \caption{Stability of the Monte Carlo dropout uncertainty estimate as a
    function of the number of stochastic forward passes $T$. The mean width of the
    $\pm 2\sigma$ band and the root mean squared error are shown for increasing
    $T$, each averaged over ten independent sets of dropout masks. The estimate
    stabilises well before the largest value of $T$, which fixes the number of
    passes used in the experiments of Chapter~\ref{ch:setup}.}
    \label{fig:mc-dropout-passes}
\end{figure}


%\section{Laplace approximation}
\section{Laplace approximation}
\label{sec:laplace-approximation}

Section~\ref{subsec:mle-map} left the MAP estimate as an unfinished result: a
single point that discards every trace of how many other parameter settings
explained the data equally well; however, the Laplace approximation recovers part of that
information without discarding the point. It fits a Gaussian centered on
$\theta_{\mathrm{MAP}}$, with a covariance matrix taken from the local curvature
of the loss, and thereby turns a trained point-estimate network into a Bayesian
one~\cite{daxberger2021}. The method was introduced to neural networks by
MacKay~\cite{mackay1992}, and the two scalable treatments used below build directly on that work~\cite{ritter2018, daxberger2021}.

The facts mentioned before make the method apart from the two preceding sections; Bayes by Backprop
replaced every weight with a pair of variational parameters and optimized a new
objective. Monte Carlo dropout demanded dropout before every weight layer together with $L_2$ regularization, and Section~\ref{subsec:mc-dropout-limitations} noted how often implementations depart from that requirement. 

On the other hand, the Laplace approximation imposes no condition on training at all. It requires no modification of the training procedure, so the uncertainty of a model already in
production can be estimated without retraining it~\cite{ritter2018}, and it can be applied to a pre-trained network obtained off the shelf~\cite{daxberger2021}. It is the only post-hoc method in this comparison, and the predictive performance of the underlying point estimate is left untouched.

Let $\mathcal{L}(\mathcal{D}; \theta)$ denote the regularised training loss, that
is, the negative logarithm of the joint distribution appearing in the numerator of Equation~\eqref{eq:bayes-rule}. Its exponential is an unnormalised posterior, and $\theta_{\mathrm{MAP}}$ is a local minimum of it. Expanding $\mathcal{L}$ to second order around that minimum gives
%
\begin{equation}
    \mathcal{L}(\mathcal{D}; \theta)
    \approx \mathcal{L}(\mathcal{D}; \theta_{\mathrm{MAP}})
      + \tfrac{1}{2}\,
        (\theta - \theta_{\mathrm{MAP}})^{\!\top}
        \bigl[
          \nabla^2_\theta \mathcal{L}(\mathcal{D}; \theta)
          \bigr|_{\theta_{\mathrm{MAP}}}
        \bigr]
        (\theta - \theta_{\mathrm{MAP}}) ,
    \label{eq:laplace-taylor}
\end{equation}
%
where the first-order term is absent because the gradient vanishes at a maximum of the posterior~\cite{ritter2018}. The right-hand side of Equation~\eqref{eq:laplace-taylor} is quadratic in $\theta$, so its exponential has Gaussian functional form and the approximate posterior follows immediately,
%
\begin{equation}
    p(\theta \mid \mathcal{D})
    \approx \mathcal{N}\!\left(\theta;\, \theta_{\mathrm{MAP}},\, \Sigma\right),
    \qquad
    \Sigma :=
    \left(
      \nabla^2_\theta \mathcal{L}(\mathcal{D}; \theta)
      \bigr|_{\theta_{\mathrm{MAP}}}
    \right)^{\!-1} .
    \label{eq:laplace-posterior}
\end{equation}
%
The mean is acquired from the training and the covariance from the inverse Hessian of the loss~\cite{jospin2022}. Equation~\eqref{eq:laplace-posterior} is the whole
method; everything that follows concerns making it computable.

% ---------------------------------------------------------------------

The Hessian inherits the structure of the loss, under the Gaussian prior
$p(\theta) = \mathcal{N}(\theta;\, 0,\, \gamma^2 I)$ that corresponds to weight
decay, it splits into a prior and a likelihood contribution,
%
\begin{equation}
    \nabla^2_\theta \mathcal{L}(\mathcal{D}; \theta)
    \bigr|_{\theta_{\mathrm{MAP}}}
    = \gamma^{-2} I
      - \sum_{n=1}^{N}
        \nabla^2_\theta \log p\bigl(y_n \mid f_\theta(x_n)\bigr)
        \bigr|_{\theta_{\mathrm{MAP}}} ,
    \label{eq:laplace-hessian}
\end{equation}
%
in which the first term is diagonal and constant, and the second carries all the
information~\cite{daxberger2021}. In this equation (especially in the second therm) two obstacles arise, and they are of different kinds. The second term of Equation~\eqref{eq:laplace-hessian} grows quadratically with the number of parameters, which Section~\ref{subsec:intractability} put at $10^5$ and above, so a naive implementation is infeasible. The exact Hessian of a network is also not guaranteed to be positive semi-definite, whereas Equation~\eqref{eq:laplace-posterior} defines a distribution only if the curvature matrix is. Both obstacles are met by replacing the Hessian with a positive semi-definite surrogate: the Fisher information matrix or the generalised Gauss--Newton matrix, which coincide for the log-likelihoods in common use and have become the default choice in deep Laplace approximations~\cite{daxberger2021}.

A surrogate of the same size is still too large, so a factorization is imposed on
top of it. The lightest option keeps only the diagonal, which reduces the inversion from cubic to linear cost in the number of parameters. The approximate Kronecker-factored curvature (KFAC) sits between the diagonal and the full matrix: it models the correlations between weights within a layer and assumes independence between layers, writing each within-layer block as a
Kronecker product of two smaller matrices~\cite{daxberger2021}. Applied to the
Laplace approximation, this produces a normal posterior matrix on top of the weight
matrix of every layer, from which sampling is cheap because only the two factors
need to be inverted~\cite{ritter2018}. The choice between these structures is
not a matter of budget alone, a diagonal covariance assigns substantial
probability mass to low-probability regions of the true posterior whenever
weights are strongly correlated~\cite{ritter2018}; MacKay had already warned
against diagonal factorisation~\cite{mackay1992}. The empirical picture
agrees: the Kronecker-factored variant produced higher uncertainty on
out-of-distribution inputs and greater robustness to adversarial attacks than
its diagonal counterpart, which behaved much like dropout~\cite{ritter2018},
and diagonal approximations performed markedly worse than KFAC in a broader
benchmark~\cite{daxberger2021}.

A second route to scalability restricts inference to part of the network: subnetwork inference treats a small subset of the parameters probabilistically and leaves the rest at their MAP values; the last layer variant, which fixes the first $L-1$ layers as a feature extractor, is the special case recommended as a default because it is cheap and, for the last layer, the Hessian coincides with the generalized Gauss--Newton matrix~\cite{daxberger2021}. Neither restriction is adopted here, the networks compared in this work contain a few thousand parameters, so the full curvature matrix remains affordable, and the covariance structure can be treated as an experimental variable rather than as a concession to cost. This also keeps the comparison honest because the remaining methods place a distribution over all weights.

The prior precision deserves separate treatment, since it is where the post-hoc character of the method becomes a practical freedom. Separation of training from inference means that the prior variance $\gamma^2$ entering Equation~\eqref{eq:laplace-hessian} does not need to be equal to the weight decay used during training~\cite{daxberger2021}. Ritter et al.\ exploit this by treating the prior precision and scale of the curvature factors as hyperparameters tuned in a validation set, which requires no retraining and is trivial to parallelise~\cite{ritter2018}. A principled alternative dispenses with the validation set and maximizes the Laplace approximation to the evidence, in the tradition of the MacKay framework~\cite{mackay1992}, in which the evidence also provides an objective choice of the weight decay term~\cite{daxberger2021}. The quantity being maximized is the same one that Section~\ref{subsec:gp-hyperparameters} obtained analytically for a Gaussian process in Equation~\eqref{eq:gp-marginal-likelihood}. For a network, it is itself
an approximation, which is a fair summary of the difference between the two methods.

% ---------------------------------------------------------------------
\subsection{Prediction and uncertainty decomposition}
\label{subsec:laplace-prediction}

A Gaussian posterior over the weights does not induce a Gaussian distribution over the outputs, because the network is non-linear. Two ways of evaluating Equation~\eqref{eq:predictive-2} follow from this, the first draws $T$ weight samples from Equation~\eqref{eq:laplace-posterior} and averages the resulting predictions~\cite{ritter2018}, which reproduces the estimator already used in Equation~\eqref{eq:bbb-variance-estimator}. The second linearises the network about the MAP estimate, $f_\theta(x^\ast) \approx f_{\theta_{\mathrm{MAP}}}(x^\ast) + J(x^\ast)^{\!\top}(\theta - \theta_{\mathrm{MAP}})$, with $J(x^\ast)$ the Jacobian of the output with respect to the parameters. The distribution over outputs is then Gaussian and available in closed form~\cite{daxberger2021}.

The choice between them is settled in the literature. Monte Carlo integration
performs poorly for Laplace approximations built on Gauss--Newton or Fisher
curvature, an inconsistency between the curvature approximation and the
predictive that the linearised predictive removes~\cite{daxberger2021}. For
regression with a Gaussian likelihood of variance $\sigma^2$, the predictive
distribution is obtained exactly by convolving the two Gaussians,
%
\begin{equation}
    p(y^\ast \mid x^\ast, \mathcal{D})
    \approx
    \mathcal{N}\!\left(
      y^\ast;\;
      f_{\theta_{\mathrm{MAP}}}(x^\ast),\;
      \underbrace{J(x^\ast)^{\!\top} \Sigma\, J(x^\ast)}_{\text{epistemic}}
      + \underbrace{\sigma^2}_{\text{aleatoric}}
    \right) .
    \label{eq:laplace-predictive}
\end{equation}
%
Equation~\eqref{eq:laplace-predictive} is the counterpart of Equation~\eqref{eq:total-variance-2} for this method, and its structure matches the Gaussian process decomposition of Equation~\eqref{eq:gp-variance-split} rather than the sampling estimators of Sections~\ref{sec:bayes-by-backprop} and~\ref{sec:mc-dropout}. The epistemic term is a quadratic form in the posterior covariance, projected onto the output through the Jacobian. It grows where the loss is flat, that is, where the data constrain the parameters weakly.

% Uwaga: człon sigma^2 w Równaniu (laplace-predictive) jest w
% daxberger2021 stałą wariancją gaussowskiej wiarygodności — czyli model
% homoskedastyczny, dokładnie jak tau^{-1} w Równaniu (mc-dropout-variance).
% Człon sigma^2 w Równaniu (laplace-predictive) jest w daxberger2021 stałą wariancją gaussowskiej %wiarygodności — model homoskedastyczny, dokładnie jak tau^{-1} w Równaniu (mc-dropout-variance). %Akapit poniżej opisuje tę własność i wskazuje warunek, przy którym człon przestaje być stały; nie %deklaruje konfiguracji eksperymentu — ta jest w Rozdziale 4.

The constant $\sigma^2$ makes Equation~\eqref{eq:laplace-predictive}
homoscedastic, exactly as the term $\tau^{-1}$ did in
Equation~\eqref{eq:mc-dropout-variance}. The restriction is not intrinsic to
the method. If the MAP model carries the variance head of
Section~\ref{subsec:aleatoric}, the likelihood variance becomes a function of
the input and $\sigma^2$ is replaced by
$\hat{\sigma}^2_{\theta_{\mathrm{MAP}}}(x^\ast)$, in the same way as in
Section~\ref{sec:mc-dropout}. Unlike there, the substitution also affects the
curvature: the Gauss--Newton matrix of Equation~\eqref{eq:laplace-hessian} is
then taken with respect to a network whose output is two-dimensional, and the
treatment of the variance head under the approximation has to be specified.
Both variants are consistent with Equation~\eqref{eq:laplace-predictive}; the
one used here is stated in Chapter~\ref{ch:setup}.

One consequence of Equation~\eqref{eq:laplace-predictive} is easy to overlook, the predictive is obtained without sampling, so a calibrated prediction costs a single forward pass, which distinguishes the method from almost all other Bayesian and ensemble approaches~\cite{daxberger2021}. The stochastic passes $T$ required by Equations~\eqref{eq:bbb-variance-estimator} and~\eqref{eq:mc-dropout-heteroscedastic} disappear entirely. The savings is not unconditional: the linearized predictive method requires a Jacobian for every test input, and that cost is the reason the technique fell out of favor for largearchitectures~\cite{daxberger2021}, on the other hand on the small networks used here the Jacobian is inexpensive, so the closed form is adopted throughout.

% [Miejsce na Rysunek 3.8: posterior Laplace na zbiorze load_sin.
%  Trening na [0, 6], ewaluacja na [-2, 10]. TE SAME osie i ta sama skala co
%  Rysunek 3.2 (GP), 3.5 (BBB) i 3.6 (MC dropout) — komplet czterech rysunków
%  na wspólnej skali jest głównym argumentem wizualnym rozdziału.
%  Sieć: MLP z src/models.py wytrenowana zwyczajnie (dropout = 0), następnie
%  Laplace policzony post hoc. Predykcyjna średnia = wyjście sieci MAP,
%  pasmo +-2 sigma z pierwiastka Równania (laplace-predictive), punkty
%  treningowe, prawdziwa funkcja linią przerywaną.
%  Cel rysunku: pokazać, że średnia jest DOKŁADNIE taka jak sieci
%  deterministycznej — czego nie ma u BBB ani MC dropout — a poza zakresem
%  treningowym pasmo rozszerza się wolniej niż u GP.]

\begin{figure}[htbp]
    \centering
    \includegraphics[width=0.85\textwidth]{rodzial3_rys/img3_8.png}
    \caption{Predictive distribution of a Laplace approximation fitted post hoc
    to a deterministically trained network on the synthetic sine dataset. The
    model is trained on the interval $[0, 6]$ and evaluated on $[-2, 10]$, with
    the axes matching those of Figures~\ref{fig:gp-posterior-sine},
    \ref{fig:bbb-posterior-sine} and~\ref{fig:mc-dropout-posterior-sine} to allow
    a direct comparison. The band shows $\pm 2\sigma$ obtained from
    Equation~\eqref{eq:laplace-predictive}. The predictive mean is that of the
    underlying point-estimate network and is therefore unchanged by the
    approximation.}
    \label{fig:laplace-posterior-sine}
\end{figure}

% ---------------------------------------------------------------------
\subsection{Computational cost and limitations}
\label{subsec:laplace-limitations}

The cost of the method is the smallest in this comparison. Training is
not affected by any additional mechanics, since the point estimate is obtained in the usual way, and the only additional step is to pass the data to accumulate the curvature factors,
performed once~\cite{ritter2018}. Measured against a MAP network, deep ensembles, variational Bayes and stochastic-gradient MCMC were between two an five times more expensive in training and prediction, whereas the Laplace approximation added negligible overhead~\cite{daxberger2021}. The memory figures tell the same story: for a wide residual network on CIFAR-10, the MAP model occupied 11~MB, the Laplace approximation 12~MB, mean-field variational Bayes 22~MB and a five-member deep ensemble 55~MB~\cite{daxberger2021}. These figures describe the recommended last-layer variant with a sampling-free predictive, so they are a lower bound rather than a general result, but the ordering is not in   doubt.

The decisive limitation is stated in the construction itself - the expansion of
Equation~\eqref{eq:laplace-taylor} is taken at one maximum, so the approximation
is local and covers a single mode of the posterior~\cite{daxberger2021}. This
matters because Section~\ref{subsec:intractability} identified the posterior of a
network as highly multimodal, with every solution replicated by the permutation
symmetries of the layers. A comparison with Bayes by Backprop sharpens the point -
both methods return a unimodal posterior, but for unrelated reasons. In Bayes by
Backprop the narrowness follows from the objective: the direction of the
divergence in Equation~\eqref{eq:vi-objective} favours mode-seeking solutions and
the mean-field family of Equation~\eqref{eq:mean-field} ignores correlations, so
a different family or a different divergence would in principle change the
outcome. In the Laplace approximation no such choice is available, the remaining
modes never enter the computation, and no covariance structure, however
expressive, can recover them. Only a mixture built from several MAP solutions
can, and both the Laplace approximation and related Gaussian approximations
admit such an extension post hoc~\cite{daxberger2021} — which anticipates the
method of Section~\ref{sec:deep-ensembles}.

A qualification is due, and it comes from the same source. Prior work and the
experiments of Daxberger et al.\ indicate that restriction to a single mode is
less limiting in practice than the argument above suggests, a behavior
attributed to the non-linear relation between parameter space and function
space~\cite{daxberger2021}. Their results support this - on in-distribution data
the Laplace approximation was the best-calibrated method in terms of expected
calibration error while retaining the accuracy of the MAP network, and it
detected out-of-distribution inputs more reliably than variational Bayes,
stochastic-gradient MCMC and SWAG~\cite{daxberger2021}. The locality of the
approximation is therefore a structural property, not a measured deficiency.

A second limitation is specific to the setting of this work: the quality of the
approximation depends on the ratio between the number of data points and the
number of curvature directions to be estimated, on a regression problem with
twenty observations and a $78 \times 78$ precision matrix, the unregularised full
Laplace approximation vastly overestimated the uncertainty and distorted the
predictive mean, because the Hessian of the log-likelihood was
underdetermined~\cite{ritter2018}. The synthetic dataset of
Section~\ref{subsec:aleatoric} places fifty training points against several
thousand parameters, so it falls squarely in that regime. Restricting the
structure of the covariance is consequently not only a computational device but
also a means of obtaining a better-conditioned estimate~\cite{ritter2018}, and
the prior precision must be tuned rather than inherited from training.

This yields an expectation for Chapter~\ref{ch:results} that differs from the one
attached to the two preceding methods. The prediction stated in
Section~\ref{subsec:bbb-limitations} and extended in
Section~\ref{subsec:mc-dropout-limitations} was based on the mode-seeking behavior
of a divergence that the Laplace approximation never minimizes, so it does not
carry over. The reported evidence points the other way: without regularisation
the intervals are too wide, and once the hyperparameters are tuned the
uncertainty is slightly higher than that of a Hamiltonian Monte Carlo reference
close to the training data and grows more slowly away from
it~\cite{ritter2018}. Two consequences follow for the metrics of
Chapter~\ref{ch:setup}: Empirical coverage may exceed the nominal level rather
than fall short of it, at the price of a larger mean interval width, and the
increase in width under extrapolation should be weaker than that of the Gaussian
process of Section~\ref{sec:gaussian-processes}. Both are read directly from
PICP, MPIW and Figure~\ref{fig:laplace-posterior-sine}.

% [Miejsce na Rysunek 3.9: wpływ struktury kowariancji na oszacowanie
%  niepewności. Zbiór load_sin, ta sama sieć MAP co na Rysunku 3.8, trzy
%  panele obok siebie na wspólnych osiach: kowariancja pełna, KFAC,
%  diagonalna. W każdym panelu pasmo +-2 sigma z Równania
%  (laplace-predictive), punkty treningowe, prawdziwa funkcja przerywaną.
%  Cel rysunku: pokazać na własnych danych to, co ritter2018 pokazuje na
%  swoim zbiorze zabawkowym — że wybór faktoryzacji zmienia szerokość
%  pasma, a wariant pełny bez regularyzacji przeszacowuje niepewność.
%  To jest jedyne miejsce w pracy, gdzie widać, że struktura kowariancji
%  jest decyzją modelową, a nie techniczną.
%  Wariant rozszerzony (jeśli starczy miejsca): druga linia paneli z tą samą
%  strukturą przy trzech wartościach prior precision, żeby uzasadnić liczbowo
%  wartość przyjętą w Rozdziale 4.]

\begin{figure}[htbp]
    \centering
    \includegraphics[width=0.95\textwidth]{rodzial3_rys/img3_9.png}
    \caption{Effect of the covariance structure on the Laplace uncertainty
    estimate, evaluated on the synthetic sine dataset with a fixed MAP network.
    The three panels show the $\pm 2\sigma$ band obtained from a full, a
    Kronecker-factored and a diagonal approximation of the curvature, on common
    axes. The choice of factorisation is a modelling decision that changes the
    reported uncertainty, and not merely a computational shortcut.}
    \label{fig:laplace-covariance-structure}
\end{figure}


%\section{Deep ensembles}
\section{Deep ensembles}
\label{sec:deep-ensembles}

The three preceding sections approximated a posterior over the weights in three
different ways: a Gaussian variational family, a Bernoulli one, and a local
Gaussian read off the curvature at a single maximum. Deep ensembles dispense
with the posterior altogether. A fixed number $M$ of networks of the same
architecture is trained independently, and their predictive distributions are
combined. Lakshminarayanan et al.\ proposed the method explicitly as an
alternative to Bayesian neural networks, on the grounds that the latter demand
substantial changes to the training pipeline and are slower to train than
standard networks~\cite{lakshminarayanan2017}. Nothing in the construction
refers to a prior, a posterior or a divergence. The method therefore enters this
comparison from the opposite direction to
Sections~\ref{sec:bayes-by-backprop}--\ref{sec:laplace-approximation}, and it
closes the chapter from the non-Bayesian side.

The motivation the authors give comes directly from
Section~\ref{sec:mc-dropout}. Dropout admits a second reading, as a combination
of an ensemble of networks with shared parameters, alongside the variational one
established by Gal and Ghahramani. Lakshminarayanan et al.\ argue that the
ensemble reading is the more plausible of the two whenever the dropout rate is
not tuned on the training data, because any sensible approximation to a
posterior must depend on that data~\cite{lakshminarayanan2017}. This is the
objection already recorded in Section~\ref{subsec:mc-dropout-limitations},
where the rate $p_i$ was seen to be chosen rather than inferred. Taken
seriously, it suggests examining ensembles on their own terms rather than as a
by-product of a variational argument.

% ---------------------------------------------------------------------
\subsection{Construction and training}
\label{subsec:ensembles-training}

The recipe consists of three elements: a proper scoring rule as the training
criterion, an ensemble, and adversarial training as an optional
addition~\cite{lakshminarayanan2017}.

A scoring rule assigns a numerical value to a predictive distribution and
rewards better calibrated predictions over worse ones. It is proper when the
true distribution attains the optimum, and only then. The log-likelihood is a
proper scoring rule, whereas the squared error of a point prediction says
nothing about predictive uncertainty. For regression, each member therefore
carries the two output heads of Section~\ref{subsec:aleatoric} and is trained by
minimising the Gaussian negative log-likelihood of
Equation~\eqref{eq:gaussian-nll}, under the heteroscedastic model of
Equation~\eqref{eq:heteroscedastic}. The variance head is passed through a
softplus transform to enforce positivity, and a small floor is added for
numerical stability~\cite{lakshminarayanan2017}.

This is a design decision rather than a formality, and the authors measured its
effect. A single network that learns its variance achieved a better test
log-likelihood than an ensemble of five, or of ten, networks trained on the
squared error and equipped with the empirical variance of their predictions. The
calibration figures make the gap concrete: on the \emph{Year Prediction MSD}
dataset, the nominal $80\%$ interval derived from that empirical variance
contained roughly $20\%$ of the test observations~\cite{lakshminarayanan2017}.
The common heuristic of ensembling point predictors and reading the spread as an
uncertainty estimate is thus not a cheap version of the method described here.
It is a different and markedly worse one.

Diversity between members is the second element, and it is obtained at no cost.
Random initialisation of the parameters, together with random shuffling of the
training points, was found sufficient in practice. Every member is trained on
the whole training set rather than on a bootstrap sample, because bagging
degraded performance in the reported experiments; a bootstrap sample contains on
average only $63.2\%$ unique points, and a base learner with many local optima
does not need the extra randomisation~\cite{lakshminarayanan2017}. The same
protocol is adopted in this work, which also keeps the training data identical
to that used by the other four methods.

Why so little randomisation suffices was examined later by Fort et
al.~\cite{fort2019}. Independently initialised networks reach entirely different
modes of the loss landscape. Functions collected along a single optimisation
trajectory, or sampled from a subspace around it, cluster within one mode in the
space of predictions, even when they differ considerably in weight space.
Low-loss paths connecting distinct modes do exist, but the functions along them
are not connected in prediction space. Measured on the diversity--accuracy
plane, the decorrelating power of random initialisation is not matched by
subspace sampling methods~\cite{fort2019}.

This is the mechanism behind the epistemic estimates of the method, and it
locates the contrast with the two preceding sections precisely. Bayes by
Backprop is unimodal because the divergence of
Equation~\eqref{eq:vi-objective} is mode-seeking and the family of
Equation~\eqref{eq:mean-field} ignores correlations, as set out in
Section~\ref{subsec:bbb-limitations}. The Laplace approximation is unimodal
because the expansion is taken at one maximum, as set out in
Section~\ref{subsec:laplace-limitations}. Fort et al.\ propose exactly this
difference as the explanation for the gap between theory and practice: scalable
variational methods concentrate on a single mode, while deep ensembles cover
several, and Bayesian neural networks underperform ensembles in practice
particularly under dataset shift~\cite{fort2019}.

The third element is adversarial training. Given an input $x$ with target $y$
and a loss $\ell(\theta, x, y)$, the fast gradient sign method produces a
perturbed example
%
\begin{equation}
    x' = x + \epsilon \operatorname{sign}
      \bigl( \nabla_{x}\, \ell(\theta, x, y) \bigr) ,
    \label{eq:fgsm}
\end{equation}
%
which is added to the training set with the original target. The perturbation
follows the direction in which the loss rises fastest, so the procedure is read
as a cheap way of smoothing the predictive distribution over a neighbourhood of
each training point, rather than along a randomly chosen
direction~\cite{lakshminarayanan2017}. The step is marked optional in the
original algorithm, and the reported evidence is mixed. On MNIST it improved
every metric at every ensemble size, on SVHN the benefit faded as $M$ grew, and
on the regression benchmarks the difference between an ensemble with and without
it lay within the error bars~\cite{lakshminarayanan2017}.

Adversarial training was excluded from this work, for two reasons. The reported
gain on regression is not distinguishable from noise, and the perturbation size
$\epsilon$ is an additional hyperparameter that would have to be set per
dataset. Tuning it for one method alone would make the comparison depend on the
effort spent on that method, which is the argument already used for the kernel
choice in Section~\ref{subsec:gp-kernel}. A second consideration is that
Equation~\eqref{eq:fgsm} alters the training set, so the ensemble would no
longer be trained on the same data as the other four methods.

% ---------------------------------------------------------------------
\subsection{Prediction and uncertainty decomposition}
\label{subsec:ensembles-prediction}

The trained members are treated as a uniformly weighted mixture,
%
\begin{equation}
    p\bigl(y^\ast \mid x^\ast\bigr)
    = \frac{1}{M} \sum_{m=1}^{M}
      p_{\theta_m}\bigl(y^\ast \mid x^\ast\bigr) ,
    \label{eq:ensemble-mixture}
\end{equation}
%
where $\theta_m$ denotes the parameters of the $m$-th network. For Gaussian
members this is a mixture of Gaussians. To keep quantiles and predictive
probabilities easy to compute, the mixture is approximated by a single Gaussian
carrying its mean and its variance~\cite{lakshminarayanan2017}.

Both moments are already available. The mean is the average of the member means,
that is Equation~\eqref{eq:bbb-predictive-mean} with the index over stochastic
passes replaced by an index over members. The variance is
Equation~\eqref{eq:bbb-variance-estimator} under the same substitution, and it
is identical in form to Equation~\eqref{eq:mc-dropout-heteroscedastic}. The
estimator is therefore the third appearance of one expression, and it is not
restated here. What differs between the three methods is the origin of the
samples, and nothing else. In Section~\ref{sec:bayes-by-backprop} they are drawn
from a learnt variational posterior, in Section~\ref{sec:mc-dropout} from a set
of Bernoulli masks, and here each sample is the endpoint of an independent
optimisation run. The samples of a deep ensemble are not drawn from any
distribution at all.

% Uwaga dla autora: lakshminarayanan2017 zapisuje wariancję mieszaniny
% w postaci (1/M) sum (sigma_m^2 + mu_m^2) - mu_*^2, czyli z dzielnikiem M,
% tak samo jak kendall2017 w Sekcji 3.3. Powyżej przyjęto dzielnik T-1
% z Równania (bbb-variance-estimator). To ta sama decyzja, o której mowa
% w notatce w rozdzial33.tex — zadeklarować ją raz w Rozdziale 4
% i stosować dla wszystkich metod.

The reading of the two terms carries over unchanged from
Equation~\eqref{eq:total-variance-2}. The first averages the noise predicted by
the individual networks and is the aleatoric component. The second measures how
far the members disagree about the mean and is the epistemic component. It
vanishes for $M = 1$, which restates the observation of
Section~\ref{subsec:mle-map} that a single point estimate reports no epistemic
uncertainty by assumption. Where the training data are dense, independently
trained networks agree and the term is small. Outside the training range the
members are constrained by nothing in common, they diverge, and the term grows.
The toy experiment of the original work shows precisely this: ensembling
improved the predictive distribution most as the test points moved away from the
observed data~\cite{lakshminarayanan2017}.

% [Miejsce na Rysunek 3.10: posterior deep ensemble na zbiorze load_sin.
%  Trening na [0, 6], ewaluacja na [-2, 10]. TE SAME osie i ta sama skala co
%  Rysunek 3.2 (GP), 3.5 (BBB), 3.6 (MC dropout) i 3.8 (Laplace) — komplet
%  pięciu rysunków na wspólnej skali domyka główny argument wizualny
%  rozdziału.
%  Sieć: M niezależnych MLP z src/models.py (dropout = 0), różne ziarna
%  inicjalizacji i różna kolejność danych, trening Gaussian NLL.
%  Predykcyjna średnia z Równania (bbb-predictive-mean) po składowych,
%  pasmo +-2 sigma z Równania (bbb-variance-estimator), punkty treningowe,
%  prawdziwa funkcja linią przerywaną. OBOWIĄZKOWO: M pojedynczych średnich
%  cienką linią — to jedyne miejsce w pracy, gdzie widać wprost rozbieżność
%  składowych poza zakresem treningowym, czyli mechanizm z fort2019.
%  Cel rysunku: pokazać, że pasmo rozchodzi się poza [0, 6] szerzej niż
%  u BBB i MC dropout, i porównać tempo tego wzrostu z GP.]

\begin{figure}[htbp]
    \centering
    \includegraphics[width=0.85\textwidth]{rodzial3_rys/img3_10.png}
    \caption{Predictive distribution of a deep ensemble on the synthetic sine
    dataset. The models are trained on the interval $[0, 6]$ and evaluated on
    $[-2, 10]$, with the axes matching those of
    Figures~\ref{fig:gp-posterior-sine}, \ref{fig:bbb-posterior-sine},
    \ref{fig:mc-dropout-posterior-sine} and~\ref{fig:laplace-posterior-sine} to
    allow a direct comparison. Thin lines show the predictive means of the
    individual members, which coincide within the training range and separate
    outside it. The band shows $\pm 2\sigma$ obtained from the mixture of
    Equation~\eqref{eq:ensemble-mixture}.}
    \label{fig:ensemble-posterior-sine}
\end{figure}

% [Miejsce na Rysunek 3.11: wpływ liczby składowych M na oszacowanie
%  niepewności. Zbiór load_sin, ta sama architektura co na Rysunku 3.10.
%  Średnia szerokość pasma +-2 sigma oraz RMSE w funkcji M = 1, 2, 3, 5, 10,
%  15, 20, każdy punkt uśredniony po kilku losowaniach zestawu ziaren,
%  z pasmem odchylenia standardowego między losowaniami. Rysunek jest
%  odpowiednikiem Rysunku 3.7 dla MC dropout i ma tę samą rolę: uzasadnić
%  liczbowo wartość M przyjętą w Rozdziale 4, zamiast przyjmować ją
%  arbitralnie. Punkt odniesienia z literatury: M = 5 jako wartość domyślna
%  u lakshminarayanan2017 i malejące przyrosty powyżej pięciu składowych
%  u ovadia2019.]

\begin{figure}[htbp]
    \centering
    \includegraphics[width=0.85\textwidth]{rodzial3_rys/img3_11.png}
    \caption{Quality of the deep ensemble uncertainty estimate as a function of
    the number of members $M$. The mean width of the $\pm 2\sigma$ band and the
    root mean squared error are shown for increasing $M$, each averaged over
    several independent sets of initialisation seeds. The gains diminish well
    before the largest value of $M$, which fixes the ensemble size used in the
    experiments of Chapter~\ref{ch:setup}.}
    \label{fig:ensemble-size}
\end{figure}

% [Miejsce na Rysunek 3.12 (opcjonalny): wariancja uczona kontra wariancja
%  empiryczna. Dwa panele na wspólnych osiach, oba na zbiorze load_sin:
%  po lewej M sieci trenowanych na MSE, z pasmem +-2 sigma liczonym wyłącznie
%  z empirycznego rozrzutu predykcji; po prawej ten sam ensemble trenowany
%  na Gaussian NLL, z pasmem z Równania (bbb-variance-estimator).
%  Cel rysunku: odtworzyć na własnych danych wynik lakshminarayanan2017
%  uzasadniający dwugłowicową architekturę — heurystyka „uśrednij predykcje
%  punktowe i weź rozrzut" wyraźnie zaniża niepewność. Jeśli argument
%  liczbowy w tekście wystarczy, rysunek można pominąć.]

% ---------------------------------------------------------------------
\paragraph{Computational cost and limitations.}
The cost of the method is the highest in this comparison, and it is also the
easiest to describe. Training is repeated $M$ times and prediction requires $M$
forward passes, so both scale linearly with the ensemble size. The saving
argument advanced in Section~\ref{subsec:mc-dropout-limitations} applies here
without the qualification attached to it there: the members are trained
independently and share nothing, so the whole procedure parallelises trivially
and the wall-clock cost need not grow at all~\cite{lakshminarayanan2017}. Memory
is the component that does not parallelise away. The ensemble holds $M$ complete
copies of the network, which is the only entry in the cost table of Ovadia et
al.\ whose storage grows with the number of replications; Monte Carlo dropout
adds nothing and mean-field variational inference doubles the
count~\cite{ovadia2019}. Where memory is the binding constraint, the ensemble
can be distilled into a single smaller model~\cite{lakshminarayanan2017}.

Two limitations follow from the construction. The diversity on which the
epistemic term depends is an empirical property of random initialisation, not a
guaranteed one; Fort et al.\ measured it rather than derived
it~\cite{fort2019}. Nothing in the procedure encodes beliefs held before the
data are seen, so the method offers no counterpart to the kernel of
Section~\ref{subsec:gp-kernel} or the prior precision of
Section~\ref{sec:laplace-approximation} through which such beliefs could be
stated or tuned. A third qualification is specific to the present work. The
strongest evidence for deep ensembles under distribution shift comes from
classification benchmarks~\cite{ovadia2019, wilson2020}, whereas the comparison
in Chapter~\ref{ch:results} is a regression one, so those conclusions are
carried over as an expectation rather than as an established result.

That evidence is nevertheless consistent. Ovadia et al.\ benchmarked the
methods of this chapter across image, text and categorical data under
increasing corruption of the inputs. Their summary is that the quality of
uncertainty degrades with shift for every method, that calibration on
independent and identically distributed test data does not carry over to
shifted data, and that deep ensembles perform best across most metrics while
proving the most robust to shift~\cite{ovadia2019}. The advantage was shown not
to be a matter of capacity, since widening the single-network baselines did not
reproduce it, and a modest ensemble of five members already captured most of the
available gain~\cite{ovadia2019}.

A closing observation returns the method to the framework of
Chapter~\ref{ch:methods}. Wilson and Izmailov argue that the distinguishing
property of a Bayesian approach is marginalisation rather than optimisation, and
that deep ensembles are not a competitor to Bayesian inference but a mechanism
for performing it~\cite{wilson2020}. Because a network is typically
underspecified by the data, many different parameter settings explain it
equally well, which is exactly the situation in which averaging over them
matters. By representing several basins of attraction, an ensemble can
approximate the integral of Equation~\eqref{eq:predictive-2} more faithfully
than a method that represents one basin accurately. Their experiment supports
the claim directly. Against a reference predictive distribution built from long
Hamiltonian Monte Carlo chains, the deep ensemble was closer to the reference
than factorised variational inference, which was strongly overconfident between
clusters of data points; the distance to the reference fell with the number of
ensemble members, while for the variational method it barely moved with the
number of samples~\cite{wilson2020}.

This is the note on which the chapter closes.
Section~\ref{subsec:predictive-distribution} stated that the five methods differ
only in how they approximate Equation~\eqref{eq:predictive-2}. The last of them
was introduced as an alternative to Bayesian inference and can be read as a
particularly effective instance of it, which suggests that the label attached to
a method orders it less usefully than the quality of the approximation it
delivers. Chapter~\ref{ch:results} is arranged to test that ordering
empirically, and the expectation attached to deep ensembles differs from those
recorded for the earlier methods. Prediction intervals should be wider than
those of Bayes by Backprop and Monte Carlo dropout, empirical coverage should
sit closer to the nominal level than for either, and the growth of the epistemic
term under extrapolation should be the strongest among the four neural methods,
with the Gaussian process of Section~\ref{sec:gaussian-processes} as the
reference. The PICP and MPIW metrics of Chapter~\ref{ch:setup}, together with
Figure~\ref{fig:ensemble-posterior-sine}, are reported to settle it.



%\chapter{Experimental setup and Methodology}
%\section{Implementation details}
%\subsection{Base network architecture}
%\subsection{Influence of Depth/Width on Uncertainty Estimates *raczej do wyboru 1 opcja*}
%\subsection{Uncertainty Method Configurations}
%\section{Datasets}
%\subsection{Synthetic problem *ten sinus albo coś podobnego czy warto dodać ?*}
%\subsection{Benchmark Datasets}
%\section{Evaluation metrics}
\include{rozdzial4}







\chapter{Results}
\section{Analysis on Synthetic Data *jeśli zostaje ten przykład*}
\section{Evaluation on Benchmark Datasets}
\section{Impact of Model Topology on Uncertainty *GŁĘBOKOŚC/SZEROKOSC*}

    
\chapter{Conclusions and future work}


\chapter{References} %% \bibliography{}


 
\end{document}